\section{Dérivations et différentielles}
\label{section:0.20-fr}
\label{0.20-fr}

Les notions introduites dans ce paragraphe seront reprises sous forme
géométrique au chapitre~IV, §~16, et joueront un rôle important dans l'étude
des préschémas. Leur importance dans ce chapitre tient d'abord à leurs
relations avec la notion d'algèbre formellement lisse, et notamment aux
théorèmes (20.4.9), (20.5.7) et (20.5.12), qui seront traduits en langage
géométrique dans le paragraphe du chapitre~IV consacré aux morphismes lisses,
et sont d'un usage constant dans les applications. D'autre part, les notions
différentielles serviront à démontrer au §~22 d'importants critères de
régularité qui joueront un rôle essentiel dans l'étude approfondie des anneaux
locaux noethériens faite au §~7 du chapitre~IV.

\subsection{Dérivations et extensions d'algèbres}
\label{subsection:0.20.1-fr}

\oldpage[0\textsubscript{IV-1}]{117}
\begin{proposition}[20.1.1]
\label{0.20.1.1-fr}
Soient $A$ un anneau (non nécessairement commutatif), $B$, $E$, $C$ des
$A$-anneaux, $p:E\to C$ un $A$-homomorphisme dont le noyau $\mathfrak{Q}$
est de carré nul, $u:B\to C$ un $A$-homomorphisme. Supposons qu'il existe
un $A$-homomorphisme $v_0:B\to E$ tel que $u$ se factorise en
$B\xrightarrow{v_0}E\xrightarrow{p}C$. Alors l'ensemble des
$A$-homomorphismes $v:B\to E$ tels que $u=p\circ v$ est identique à
l'ensemble des applications $v_0+D$, où $D:B\to\mathfrak{Q}$ est une
application vérifiant les deux conditions suivantes :
\begin{enumerate}
  \item[\rm(i)] $D$ est un homomorphisme de $A$-bimodules.
  \item[\rm(ii)] Quels que soient $f$, $g$ dans $B$, on a
\end{enumerate}
\begin{equation}
\tag{20.1.1.1}
\label{0.20.1.1.1-fr}
  D(fg)=f.D(g)+D(f).g.
\end{equation}
\end{proposition}

\begin{proof}
Dire que $p(v(f))=p(v_0(f))$ pour $f\in B$ signifie que
$D(f)=v(f)-v_0(f)$ appartient à $\mathfrak{Q}$ ; si l'on écrit que
$v(fg)=v(f)v(g)$, on obtient la relation \sref{0.20.1.1.1-fr},
$\mathfrak{Q}$ étant de carré nul, et la condition (i) résulte de ce que $v$
et $v_0$ sont des $A$-homomorphismes.

Si $\rho:A\to B$ est l'homomorphisme structural, on tire de
\sref{0.20.1.1.1-fr} que l'on a
$D(\rho(a)f)=\rho(a)D(f)+D(\rho(a))f$ pour tout $a\in A$ ; mais on doit
avoir $D(\rho(a)f)=\rho(a)D(f)$ d'après (i), donc, en faisant $f=1$, il
vient
\[
  D(\rho(a))=0 \qquad\text{pour}\quad a\in A\,;
\]
réciproquement, si $D$ est nulle dans $\rho(A)$ et vérifie
\sref{0.20.1.1.1-fr}, elle vérifie aussi (i).
\end{proof}

\begin{definition}[20.1.2]
\label{0.20.1.2-fr}
Étant donnés un $A$-anneau $B$ et un $B$-bimodule $L$, on appelle
\emph{$A$-dérivation de $B$ dans $L$} une application $D:B\to L$ vérifiant
les conditions (i) et (ii) de \sref{0.20.1.1-fr}.
\end{definition}

Il résulte de \sref{0.20.1.1-fr} que le \emph{noyau} d'une $A$-dérivation
$D:B\to L$ est un sous-$A$-anneau de $B$.

On appelle parfois \emph{dérivation} de $B$ dans $L$ une
$\mathbf{Z}$-dérivation, c'est-à-dire une application \emph{additive} de $B$
dans $L$ vérifiant \sref{0.20.1.1.1-fr} ; on a donc $D(1)=0$ pour une telle
application. Si $B$ est une algèbre sur un \emph{corps premier} $P$, toute
$\mathbf{Z}$-dérivation de $B$ est une $P$-dérivation : c'est évident d'après
ce qui précède si $P$ est de caractéristique $>0$, et dans le cas contraire,
la relation $n.n^{-1}=1$ pour $n\in\mathbf{Z}$ donne
\[
  nD(n^{-1})=0
\]
donc $D(n^{-1})=0$ puisque $P$ est de caractéristique $0$.

Il résulte aussitôt de la définition \sref{0.20.1.2-fr} que si $D$, $D'$ sont
deux $A$-dérivations de $B$ dans $L$, il en est de même de $D-D'$. Autrement
dit, l'ensemble des $A$-dérivations de $B$ dans $L$ est muni d'une structure
de \emph{groupe additif} ; on note ce groupe
$\operatorname{D\acute{e}r}_A(B,L)$. Si $A$ est \emph{commutatif}, $B$ une
$A$-\emph{algèbre commutative}, et $L$ un $B$-module, alors, pour toute
$A$-dérivation $D$ de $B$ dans $L$ et tout $a\in A$, $aD$ est encore une
$A$-dérivation de $B$ dans $L$, autrement dit $\operatorname{D\acute{e}r}_A(B,L)$
est alors muni d'une structure de $A$-\emph{module}.

\oldpage[0\textsubscript{IV-1}]{118}
La prop. \sref{0.20.1.1-fr} s'interprète de la façon suivante :

\begin{corollary}[20.1.3]
\label{0.20.1.3-fr}
Étant donnés deux $A$-homomorphismes de $A$-anneaux $p:E\to C$,
$u:B\to C$ dont le premier a un noyau $\mathfrak{Q}$ de carré nul,
l'ensemble des $A$-homomorphismes $v:B\to E$ tels que $u=p\circ v$ est vide
ou est un espace homogène principal pour le groupe
$\operatorname{D\acute{e}r}_A(B,\mathfrak{Q})$.
\end{corollary}

En particulier :

\begin{corollary}[20.1.4]
\label{0.20.1.4-fr}
Soient $A$ un anneau, $C$ un $A$-anneau, $L$ un $C$-bimodule, $E$ une
$A$-extension de $C$ par $L$, $p:E\to C$ l'augmentation. L'application qui,
à toute dérivation $D\in\operatorname{D\acute{e}r}_A(C,L)$, associe l'application
$x\mapsto x+D(p(x))$ est un isomorphisme du groupe
$\operatorname{D\acute{e}r}_A(C,L)$ sur le groupe des $A$-équivalences de $E$ sur
elle-même.
\end{corollary}

\begin{proof}
Appliquons \sref{0.20.1.1-fr} pour $B=E$, $v_0=\mathrm{I}_E$ : l'ensemble
des $A$-homomorphismes $v:E\to E$ tels que $p\circ v=p$ est identique à
l'ensemble des applications $v=\mathrm{I}_E+D'$, où
$D'\in\operatorname{D\acute{e}r}_A(E,L)$. Dire qu'un tel $A$-homomorphisme $v$ est
une $A$-équivalence équivaut à dire que $v$ se réduit à l'injection canonique
dans $L$, ou encore que $D'(x)=0$ dans $L$ ; mais cela signifie aussi que
$D'$ se factorise en $E\xrightarrow{p}C\xrightarrow{D}L$, où $D$ est une
$A$-dérivation. D'où le corollaire.
\end{proof}

Pour les extensions triviales, \sref{0.20.1.3-fr} donne :

\begin{corollary}[20.1.5]
\label{0.20.1.5-fr}
Soient $A$ un anneau, $B$, $C$ deux $A$-anneaux, $L$ un $C$-bimodule,
$u:B\to C$ un $A$-homomorphisme ; l'application qui, à toute dérivation
$D\in\operatorname{D\acute{e}r}_A(B,L)$, associe l'application
$\alpha:x\mapsto(u(x),D(x))$ est une bijection sur l'ensemble des
$A$-homomorphismes $v:B\to D_C(L)$ (cf. 18.2.3) tels que $u$ se factorise en
$B\xrightarrow{v}D_C(L)\to C$.
\end{corollary}

Plus particulièrement :

\begin{corollary}[20.1.6]
\label{0.20.1.6-fr}
Soient $A$ un anneau, $B$ un $A$-anneau, $L$ un $B$-bimodule. Si, à toute
dérivation $D\in\operatorname{D\acute{e}r}_A(B,L)$, on associe : $1^\circ$ la
$A$-équivalence $(x,y)\mapsto(x,y+D(x))$ de l'extension $D_B(L)$ sur
elle-même ; $2^\circ$ le $A$-homomorphisme $x\mapsto(x,D(x))$ de $B$ dans
$D_B(L)$, inverse à droite de l'homomorphisme d'augmentation $D_B(L)\to B$,
on définit des correspondances biunivoques canoniques entre :
\begin{enumerate}
  \item[\rm(i)] l'ensemble $\operatorname{D\acute{e}r}_A(B,L)$ ;
  \item[\rm(ii)] l'ensemble des $A$-équivalences de $D_B(L)$ sur elle-même ;
  \item[\rm(iii)] l'ensemble des $A$-homomorphismes inverses à droite de
  l'homomorphisme d'augmentation $D_B(L)\to B$.
\end{enumerate}
\end{corollary}

On notera que la correspondance biunivoque ainsi établie entre (i) et (ii)
respecte les structures de groupe, et que celle qu'on en déduit entre (ii) et
(iii) n'est autre que la correspondance biunivoque déjà définie dans (18.3.8).

\begin{corollary}[20.1.7]
\label{0.20.1.7-fr}
Soient $A$ un anneau, $C$ un $A$-anneau, $L$ un $A$-bimodule, $E$ une
$A$-extension triviale de $C$ par $L$, $p:E\to C$ l'augmentation. L'ensemble
des $A$-dérivations $D$ de $E$ dans $L$ telles que $D(x)=x$ dans $L$ est
identique à l'ensemble des applications $x\mapsto x-w(p(x))$, où $w$
parcourt l'ensemble des $A$-homomorphismes inverses à droite de $p$.
\end{corollary}

\begin{proof}
Il suffit encore d'appliquer \sref{0.20.1.1-fr} pour $B=E$,
$v_0=\mathrm{I}_E$ ; si $v=\mathrm{I}_E-D$, la condition $D(x)=x$ pour
$x\in L$ équivaut à $v(x)=0$ pour $x\in L$, c'est-à-dire à $v=w\circ p$,
où $w:C\to E$ est un $A$-homomorphisme ; en outre la condition
$p\circ v=p$ équivaut à $p\circ w=\mathrm{I}_C$, autrement dit au fait que
$w$ est inverse à droite de $p$.
\end{proof}

\subsection{Propriétés fonctorielles des dérivations}
\label{subsection:0.20.2-fr}

\oldpage[0\textsubscript{IV-1}]{119}
\begin{env}[20.2.1]
\label{0.20.2.1-fr}
Soient $A$ un anneau, $B$ un $A$-anneau, $L$ un $B$-bimodule ; si $L'$ est
un second $B$-bimodule et $w:L\to L'$ un homomorphisme de $B$-bimodules, il
est clair que l'application $D\mapsto w\circ D$ est un homomorphisme de
groupes additifs
\begin{equation}
\tag{20.2.1.1}
\label{0.20.2.1.1-fr}
  w_0:\operatorname{D\acute{e}r}_A(B,L)\to\operatorname{D\acute{e}r}_A(B,L')
\end{equation}
et que si $w':L'\to L''$ est un second homomorphisme de $B$-bimodules, on a
$(w'\circ w)_0=w'_0\circ w_0$. Lorsque $A$ est commutatif, $B$ une
$A$-algèbre commutative et $L$ un $B$-module, \sref{0.20.2.1.1-fr} est un
homomorphisme de $A$-modules.

En second lieu, soient $B'$ un $A$-anneau, $v:B'\to B$ un
$A$-homomorphisme qui fait de $L$ un $B'$-bimodule ; alors l'application
$D\mapsto D\circ v$ est un homomorphisme de groupes additifs,
\begin{equation}
\tag{20.2.1.2}
\label{0.20.2.1.2-fr}
  v^0:\operatorname{D\acute{e}r}_A(B,L)\to\operatorname{D\acute{e}r}_A(B',L)
\end{equation}
comme il résulte de \sref{0.20.1.1.1-fr} ; si $v':B''\to B'$ est un second
$A$-homomorphisme, on a $(v\circ v')^0=v'^0\circ v^0$. Lorsque $A$, $B$ et
$B'$ sont commutatifs et $L$ un $B$-module, \sref{0.20.2.1.2-fr} est un
homomorphisme de $A$-modules.

Enfin, soit $u:A'\to A$ un homomorphisme d'anneaux faisant de $B$ un
$A'$-anneau ; toute $A$-dérivation $D\in\operatorname{D\acute{e}r}_A(B,L)$ est aussi
une $A'$-dérivation, d'où une injection canonique de groupes commutatifs
\begin{equation}
\tag{20.2.1.3}
\label{0.20.2.1.3-fr}
  u^0:\operatorname{D\acute{e}r}_A(B,L)\to\operatorname{D\acute{e}r}_{A'}(B,L)
\end{equation}
et si $u':A''\to A'$ est un second homomorphisme d'anneaux, on a
$(u\circ u')^0=u'^0\circ u^0$ ; lorsque $A$, $A'$ et $B$ sont commutatifs et
$L$ un $B$-module, \sref{0.20.2.1.3-fr} est un di-homomorphisme de modules
(relatif à $u$).

On peut encore dire que
\[
  (A,B,L)\leadsto\operatorname{D\acute{e}r}_A(B,L)
\]
est un \emph{foncteur covariant} de la catégorie $\boldsymbol{K}$ définie dans
(18.3.5) dans la catégorie $\boldsymbol{Ab}$ des groupes commutatifs, en
faisant correspondre, à tout triplet $(u,v,w)$ constituant un morphisme de
$\boldsymbol{K}$ l'homomorphisme $w_0\circ v^0\circ u^0$ ; la vérification de
la fonctorialité résulte de la commutativité des diagrammes
\[
\xymatrix{
  \operatorname{D\acute{e}r}_A(B,L) \ar[r]^{v^0} \ar[d]_{w_0} &
    \operatorname{D\acute{e}r}_A(B',L) \ar[d]^{w_0} \\
  \operatorname{D\acute{e}r}_A(B,L') \ar[r]^{v^0} &
    \operatorname{D\acute{e}r}_A(B',L')
}
\qquad
\xymatrix{
  \operatorname{D\acute{e}r}_A(B,L) \ar[r]^{u^0} \ar[d]_{w_0} &
    \operatorname{D\acute{e}r}_{A'}(B,L) \ar[d]^{w_0} \\
  \operatorname{D\acute{e}r}_A(B,L') \ar[r]^{u^0} &
    \operatorname{D\acute{e}r}_{A'}(B,L')
}
\]
pour tout homomorphisme $w:L\to L'$ de $B$-bimodules.
\end{env}

\oldpage[0\textsubscript{IV-1}]{120}
\begin{theorem}[20.2.2]
\label{0.20.2.2-fr}
Soient $u:A\to B$, $v:B\to C$ deux homomorphismes d'anneaux, $L$ un
$C$-bimodule. On a une suite exacte canonique de groupes commutatifs
\begin{align}
\tag{20.2.2.1}
\label{0.20.2.2.1-fr}
  0\to\operatorname{D\acute{e}r}_B(C,L)\xrightarrow{u^0}
  \operatorname{D\acute{e}r}_A(C,L)&\xrightarrow{v^0}
  \operatorname{D\acute{e}r}_A(B,L)\xrightarrow{\partial} \\
  &\to\operatorname{Exan}_B(C,L)\xrightarrow{u^1}
  \operatorname{Exan}_A(C,L)\xrightarrow{v^1}
  \operatorname{Exan}_A(B,L) \notag
\end{align}
où $u^0$, $v^0$ sont les homomorphismes (20.2.1.3) et (20.2.1.2)
respectivement, $u^1$, $v^1$ les homomorphismes définis dans (18.3.4.1) et
(18.3.3.1) respectivement, et où $\partial$ est défini comme suit : pour
toute $A$-dérivation $D$ de $B$ dans $L$, $\partial(D)$ est la classe de la
$B$-extension de $C$ par $L$ définie sur l'anneau $D_C(L)$ par le
$A$-homomorphisme $\alpha:x\mapsto(v(x),D(x))$ (cf.
\sref{0.20.1.5-fr}). En outre, la suite exacte \sref{0.20.2.2.1-fr} est
fonctorielle en $L$ (pour les homomorphismes définis dans (20.2.1.1) et
(18.3.1.1) respectivement).
\end{theorem}

\begin{proof}
Comme $D$ est une $A$-dérivation (et \emph{a fortiori} une
$\mathbf{Z}$-dérivation) de $B$ dans $L$, le $A$-homomorphisme
$\alpha:x\mapsto(v(x),D(x))$ définit bien sur $D_C(L)$ une structure de
$B$-extension, donc $\partial(D)$ est bien définie \sref{0.20.1.5-fr}. Il
faut vérifier l'exactitude en $5$ endroits :

\medskip
\noindent 1)\enspace
L'exactitude en $\operatorname{D\acute{e}r}_B(C,L)$ est triviale (cf.
\sref{0.20.2.1-fr}).

\medskip
\noindent 2)\enspace
Par définition \sref{0.20.2.1-fr} le noyau de $v^0$ est l'ensemble des
$A$-dérivations de $C$ dans $L$ qui s'annulent dans $v(B)$, c'est-à-dire
celles des $A$-dérivations qui sont aussi des $B$-dérivations
\sref{0.20.1.1-fr} ; d'où l'exactitude en $\operatorname{D\acute{e}r}_A(C,L)$.

\medskip
\noindent 3)\enspace
Le noyau de $\partial$ est formé des dérivations
$D\in\operatorname{D\acute{e}r}_A(B,L)$ pour lesquelles la $B$-extension définie par
$\alpha:x\mapsto(v(x),D(x))$ est $B$-triviale ; cela signifie (18.2.3) qu'il
existe un $B$-homomorphisme $z\mapsto(z,w(z))$ de $C$ dans $D_C(L)$ (la
structure de $B$-anneau de $D_C(L)$ étant définie par $\alpha$) ; mais un tel
homomorphisme, étant \emph{a fortiori} un $A$-homomorphisme, est de la forme
$z\mapsto(z,D'(z))$ où $D'\in\operatorname{D\acute{e}r}_A(C,L)$
\sref{0.20.1.6-fr} ; et en écrivant que c'est un $B$-homomorphisme, il vient
\[
  D(x)z+v(x)D'(z)=D'(v(x)z)=v(x)D'(z)+D'(v(x))z
\]
pour $x\in B$, $z\in C$, ce qui donne $D'(v(x))=D(x)$ ; le noyau de
$\partial$ est donc l'image de $v^0$.

\medskip
\noindent 4)\enspace
Le noyau de $u^1$ est l'ensemble des classes de $B$-extensions de $C$ par
$L$ qui sont $A$-triviales (18.3.7), donc (à équivalence près) de la forme
$D_C(L)$, où la structure de $A$-anneau est définie par l'homomorphisme
$t\mapsto(u(t),0)$. Or, toute structure de $B$-extension sur $D_C(L)$ est
définie par un homomorphisme $\alpha:x\mapsto(v(x),D(x))$, où $D$ est une
$\mathbf{Z}$-dérivation de $B$ dans $L$ \sref{0.20.1.5-fr} ; dire que la
structure de $A$-anneau de cette $B$-extension est déduite de sa structure de
$B$-anneau au moyen de $u:A\to B$ signifie que $D(u(t))=0$ pour $t\in A$,
donc que $D$ est une $A$-dérivation ou encore que la classe de la $B$-extension
considérée est de la forme $\partial(D)$ ; d'où l'exactitude en
$\operatorname{Exan}_B(C,L)$.

\medskip
\noindent 5)\enspace
Le noyau de $v^1$ est l'ensemble des classes des $A$-extensions $E$ de $C$
par $L$ qui se trivialisent sur $v(B)$, c'est-à-dire pour lesquelles il y a un
$A$-homomorphisme $w:B\to E$ tel que $v$ se factorise en
$B\xrightarrow{w}E\to C$ ; or un tel $A$-homomorphisme définit sur $E$ une
structure de $B$-extension dont la classe a pour image par $u^1$ la classe de
la $A$-extension donnée ; la réciproque étant triviale, on a prouvé
l'exactitude en $\operatorname{Exan}_A(C,L)$.

Enfin, la fonctorialité en $L$ découle trivialement des définitions.
\end{proof}

\oldpage[0\textsubscript{IV-1}]{121}
\begin{corollary}[20.2.3]
\label{0.20.2.3-fr}
Soient $A$, $B$, $C$ trois anneaux commutatifs, $u:A\to B$, $v:B\to C$
deux homomorphismes d'anneaux, $L$ un $C$-module. On a une suite exacte
canonique de $A$-modules
\begin{align}
\tag{20.2.3.1}
\label{0.20.2.3.1-fr}
  0\to\operatorname{D\acute{e}r}_B(C,L)\xrightarrow{u^0}
  \operatorname{D\acute{e}r}_A(C,L)&\xrightarrow{v^0}
  \operatorname{D\acute{e}r}_A(B,L)\xrightarrow{\partial} \\
  &\to\operatorname{Exalcom}_B(C,L)\xrightarrow{u^1}
  \operatorname{Exalcom}_A(C,L)\xrightarrow{v^1}
  \operatorname{Exalcom}_A(B,L) \notag
\end{align}
fonctorielle en $L$.
\end{corollary}

\begin{proof}
Le raisonnement est le même que dans \sref{0.20.2.2-fr} une fois qu'on a
vérifié que pour toute dérivation $D\in\operatorname{D\acute{e}r}_A(B,L)$,
$\partial(D)$ est bien la classe d'une $B$-extension de $C$ par $L$ qui est
une $B$-algèbre commutative ; mais cela résulte aussitôt de la commutativité
de $C$ et du fait que $L$ est un $C$-module.
\end{proof}

\begin{corollary}[20.2.4]
\label{0.20.2.4-fr}
Sous les hypothèses de \sref{0.20.2.2-fr} (resp.
\sref{0.20.2.3-fr}), on a une suite exacte canonique, fonctorielle en $L$,
\begin{equation}
\tag{20.2.4.1}
\label{0.20.2.4.1-fr}
  0\to\operatorname{D\acute{e}r}_B(C,L)\xrightarrow{u^0}
  \operatorname{D\acute{e}r}_A(C,L)\xrightarrow{v^0}
  \operatorname{D\acute{e}r}_A(B,L)\xrightarrow{\partial}
  \operatorname{Exan}_{B/A}(C,L)\to0
\end{equation}
(resp.
\begin{equation}
\tag{20.2.4.2}
\label{0.20.2.4.2-fr}
  0\to\operatorname{D\acute{e}r}_B(C,L)\xrightarrow{u^0}
  \operatorname{D\acute{e}r}_A(C,L)\xrightarrow{v^0}
  \operatorname{D\acute{e}r}_A(B,L)\xrightarrow{\partial}
  \operatorname{Exalcom}_{B/A}(C,L)\to0\text{).}
\end{equation}
Cela résulte de la définition de $\operatorname{Exan}_{B/A}(C,L)$ (resp.
$\operatorname{Exalcom}_{B/A}(C,L)$) ((18.3.7) et (18.4.2)).
\end{corollary}

\begin{remark}[20.2.5]
\label{0.20.2.5-fr}
Supposons que l'on ait un diagramme commutatif d'homomorphismes d'anneaux
\[
\xymatrix{
  A \ar[r] & B \ar[r] & C \\
  A' \ar[u] \ar[r] & B' \ar[u] \ar[r] & C' \ar[u]
}
\]
Alors on a un diagramme commutatif
\[
\xymatrix@C=0.5em{
  0 \ar[r] & \operatorname{D\acute{e}r}_B(C,L) \ar[r] \ar[d] &
  \operatorname{D\acute{e}r}_A(C,L) \ar[r] \ar[d] &
  \operatorname{D\acute{e}r}_A(B,L) \ar[r]^-\partial \ar[d] &
  \operatorname{Exan}_B(C,L) \ar[r] \ar[d] &
  \operatorname{Exan}_A(C,L) \ar[r] \ar[d] &
  \operatorname{Exan}_A(B,L) \ar[d] \\
  0 \ar[r] & \operatorname{D\acute{e}r}_{B'}(C',L) \ar[r] &
  \operatorname{D\acute{e}r}_{A'}(C',L) \ar[r] &
  \operatorname{D\acute{e}r}_{A'}(B',L) \ar[r]^-\partial &
  \operatorname{Exan}_{B'}(C',L) \ar[r] &
  \operatorname{Exan}_{A'}(C',L) \ar[r] &
  \operatorname{Exan}_{A'}(B',L)
}
\]
et de même pour les suites exactes \sref{0.20.2.3.1-fr},
\sref{0.20.2.4.1-fr} et \sref{0.20.2.4.2-fr}.
\end{remark}

\subsection{Dérivations continues dans les anneaux topologiques}
\label{subsection:0.20.3-fr}

\begin{env}[20.3.1]
\label{0.20.3.1-fr}
Étant donnés deux anneaux topologiques $A$, $B$ (linéairement topologisés
comme toujours), nous désignerons par $\operatorname{Hom.\ cont}(A,B)$
l'ensemble des homomorphismes continus de $A$ dans $B$. Étant donné un anneau
topologique $A$, la catégorie des $A$-anneaux topologiques se définit comme
celle des $A$-anneaux (18.1.4) en remplaçant partout « anneau » par « anneau
topologique » et « homomorphisme » par « homomorphisme continu » ; si $B$

\oldpage[0\textsubscript{IV-1}]{122}
et $C$ sont deux $A$-anneaux topologiques, on notera
$\operatorname{Hom.\ cont}_A(B,C)$ l'ensemble des $A$-homomorphismes continus
de $B$ dans $C$.

Soient $A$ un anneau topologique, $B$ un $A$-anneau topologique, $L$ un
$B$-bimodule topologique ; on note $\operatorname{D\acute{e}r.\ cont}_A(B,L)$
l'ensemble des $A$-dérivations \emph{continues} de $B$ dans $L$ ; il est clair
que c'est un \emph{sous-groupe} de $\operatorname{D\acute{e}r}_A(B,L)$ (et un
sous-$A$-module lorsque $A$ et $B$ sont commutatifs et $L$ un $B$-module).
\end{env}

\begin{env}[20.3.2]
\label{0.20.3.2-fr}
Il est immédiat que la prop. \sref{0.20.1.1-fr} subsiste quand on y remplace
« anneau » par « anneau topologique » et « homomorphisme » par
« homomorphisme continu ». De même, \sref{0.20.1.3-fr} et
\sref{0.20.1.4-fr} restent valables en remplaçant $\operatorname{D\acute{e}r}$ par
$\operatorname{D\acute{e}r.\ cont}$, « anneau » (resp. « bimodule ») par « anneau
topologique » (resp. « bimodule topologique »), « homomorphisme » par
« homomorphisme continu » ; il faut naturellement supposer dans
\sref{0.20.1.4-fr} que $p$ est continu, et remplacer « $A$-équivalences » par
« $A$-équivalences continues ». On a les mêmes résultats pour
\sref{0.20.1.5-fr} et \sref{0.20.1.6-fr}, à condition de prendre comme
topologie sur $D_C(L)$ (resp. $D_B(L)$) la topologie \emph{produit} (sur
$C\times L$, resp. $B\times L$).
\end{env}

\begin{proposition}[20.3.3]
\label{0.20.3.3-fr}
Soient $A$ un anneau topologique, $B$ un $A$-anneau topologique, $L$ un
$B$-bimodule topologique discret annulé par un idéal bilatère ouvert de $B$.
Si dans $B$ le carré de tout idéal bilatère ouvert est ouvert, on a
$\operatorname{D\acute{e}r.\ cont}_A(B,L)=\operatorname{D\acute{e}r}_A(B,L)$.
\end{proposition}

\begin{proof}
En effet, si $\mathfrak{R}$ est un idéal bilatère ouvert de $B$ annulant $L$,
et $D$ une $A$-dérivation de $B$ dans $L$, on a
$D(\mathfrak{R}^2)\subset\mathfrak{R}.D(\mathfrak{R})+
D(\mathfrak{R}).\mathfrak{R}=0$ \sref{0.20.1.1.1-fr}, donc $D$ est continue.
\end{proof}

\begin{env}[20.3.4]
\label{0.20.3.4-fr}
Tous les résultats de \sref{0.20.2.1-fr} restent valables lorsqu'on y remplace
« anneau » par « anneau topologique », « bimodule » par « bimodule
topologique », « homomorphisme » par « homomorphisme continu » et
$\operatorname{D\acute{e}r}$ par $\operatorname{D\acute{e}r.\ cont}$.
\end{env}

\begin{proposition}[20.3.5]
\label{0.20.3.5-fr}
Soient $A$ un anneau topologique, $B$ un $A$-anneau topologique, $L$ un
$B$-bimodule topologique discret annulé par un idéal bilatère ouvert de $B$.
On a alors un isomorphisme canonique
\begin{equation}
\tag{20.3.5.1}
\label{0.20.3.5.1-fr}
  \varinjlim\operatorname{D\acute{e}r}_{A/\mathfrak{J}}
  (B/\mathfrak{R},L)\xrightarrow{\sim}
  \operatorname{D\acute{e}r.\ cont}_A(B,L)
\end{equation}
où dans le premier membre la limite inductive est prise suivant l'ensemble
ordonné filtrant des couples $(\mathfrak{J},\mathfrak{R})$ d'idéaux bilatères
tels que $\mathfrak{R}.L=L.\mathfrak{R}=0$,
$\mathfrak{J}.B\subset\mathfrak{R}$,
$B.\mathfrak{J}\subset\mathfrak{R}$.
\end{proposition}

\begin{proof}
Comme $A/\mathfrak{J}$ et $B/\mathfrak{R}$ sont discrets, on a des
homomorphismes canoniques
$w_{\mathfrak{J},\mathfrak{R}}:
\operatorname{D\acute{e}r}_{A/\mathfrak{J}}(B/\mathfrak{R},L)\to
\operatorname{D\acute{e}r.\ cont}_A(B,L)$ formant un système inductif
\sref{0.20.3.4-fr}, d'où l'homomorphisme \sref{0.20.3.5.1-fr} par passage à
la limite inductive. Comme l'homomorphisme
$B/\mathfrak{R}'\to B/\mathfrak{R}$ est surjectif pour
$\mathfrak{R}\supset\mathfrak{R}'$, il résulte aussitôt de la définition que
l'homomorphisme
$\operatorname{D\acute{e}r}_{A/\mathfrak{J}}(B/\mathfrak{R},L)\to
\operatorname{D\acute{e}r}_{A/\mathfrak{J}}(B/\mathfrak{R}',L)$ (avec
$\mathfrak{R}.L=L.\mathfrak{R}=0$,
$\mathfrak{J}.B\subset\mathfrak{R}'$,
$B.\mathfrak{J}\subset\mathfrak{R}'$) est injectif, et il en est évidemment
de même de l'homomorphisme
$\operatorname{D\acute{e}r}_{A/\mathfrak{J}}(B/\mathfrak{R},L)\to
\operatorname{D\acute{e}r}_{A/\mathfrak{J}'}(B/\mathfrak{R},L)$ pour
$\mathfrak{J}'\subset\mathfrak{J}$ \sref{0.20.2.1-fr} ; on en conclut que
l'homomorphisme \sref{0.20.3.5.1-fr} est injectif. D'autre part, si $D$ est
une $A$-dérivation continue de $B$ dans $L$, son noyau contient un idéal
bilatère ouvert $\mathfrak{R}_0$ de $B$, et si $\mathfrak{J}_0$ est un idéal
bilatère ouvert de $A$ tel que $\mathfrak{J}_0B\subset\mathfrak{R}_0$ et
$B\mathfrak{J}_0\subset\mathfrak{R}_0$, il est clair que $D$ est l'image
canonique d'une $(A/\mathfrak{J}_0)$-dérivation de
$B/\mathfrak{R}_0$ dans $L$, donc \sref{0.20.3.5.1-fr} est surjectif.
\end{proof}

\oldpage[0\textsubscript{IV-1}]{123}
\begin{proposition}[20.3.6]
\label{0.20.3.6-fr}
Soient $u:A\to B$, $v:B\to C$ deux homomorphismes continus d'anneaux
topologiques, $L$ un $C$-bimodule discret annulé par un idéal bilatère ouvert
de $C$. On a une suite exacte canonique
\begin{align}
\tag{20.3.6.1}
\label{0.20.3.6.1-fr}
  0\to\operatorname{D\acute{e}r.\ cont}_B(C,L)\xrightarrow{u^0}
  \operatorname{D\acute{e}r.\ cont}_A(C,L)&\xrightarrow{v^0}
  \operatorname{D\acute{e}r.\ cont}_A(B,L)\xrightarrow{\partial} \\
  &\to\operatorname{Exantop}_B(C,L)\xrightarrow{u^1}
  \operatorname{Exantop}_A(C,L)\xrightarrow{v^1}
  \operatorname{Exantop}_A(B,L) \notag
\end{align}
où $\partial$ est défini par passage à la limite inductive à partir de
l'homomorphisme $\partial$ de \sref{0.20.2.2.1-fr} ; cette suite exacte est
fonctorielle en $L$ (dans la catégorie des $C$-bimodules discrets annulés par
des idéaux bilatères ouverts).
\end{proposition}

\begin{proof}
Cela résulte de l'exactitude du foncteur $\varinjlim$, à partir de
\sref{0.20.2.2-fr}.
\end{proof}

\begin{corollary}[20.3.7]
\label{0.20.3.7-fr}
Soient $A$, $B$, $C$ trois anneaux topologiques commutatifs, $u:A\to B$,
$v:B\to C$ deux homomorphismes continus, $L$ un $C$-module discret annulé
par un idéal ouvert de $C$. On a une suite exacte canonique de $A$-modules,
fonctorielle en $L$,
\begin{align}
\tag{20.3.7.1}
\label{0.20.3.7.1-fr}
  0\to\operatorname{D\acute{e}r.\ cont}_B(C,L)\xrightarrow{u^0}
  \operatorname{D\acute{e}r.\ cont}_A(C,L)&\xrightarrow{v^0}
  \operatorname{D\acute{e}r.\ cont}_A(B,L)\xrightarrow{\partial} \\
  &\to\operatorname{Exalcotop}_B(C,L)\xrightarrow{u^1}
  \operatorname{Exalcotop}_A(C,L)\xrightarrow{v^1}
  \operatorname{Exalcotop}_A(B,L). \notag
\end{align}
\end{corollary}

\begin{corollary}[20.3.8]
\label{0.20.3.8-fr}
Sous les hypothèses de \sref{0.20.3.5-fr} (resp.
\sref{0.20.3.6-fr}) on a une suite exacte canonique, fonctorielle en $L$
\begin{equation}
\tag{20.3.8.1}
\label{0.20.3.8.1-fr}
  0\to\operatorname{D\acute{e}r.\ cont}_B(C,L)\xrightarrow{u^0}
  \operatorname{D\acute{e}r.\ cont}_A(C,L)\xrightarrow{v^0}
  \operatorname{D\acute{e}r.\ cont}_A(B,L)\to
  \operatorname{Exantop}_{B/A}(C,L)\to0
\end{equation}
(resp.
\begin{equation}
\tag{20.3.8.2}
\label{0.20.3.8.2-fr}
  0\to\operatorname{D\acute{e}r.\ cont}_B(C,L)\xrightarrow{u^0}
  \operatorname{D\acute{e}r.\ cont}_A(C,L)\xrightarrow{v^0}
  \operatorname{D\acute{e}r.\ cont}_A(B,L)\to
  \operatorname{Exalcotop}_{B/A}(C,L)\to0\text{).}
\end{equation}
\end{corollary}

Nous laissons au lecteur le soin d'écrire les diagrammes analogues à ceux de
\sref{0.20.2.5-fr}.

\subsection{Parties principales et différentielles}
\label{subsection:0.20.4-fr}

Dans toute la suite de ce paragraphe et dans les trois suivants, tous les
anneaux sont supposés \emph{commutatifs}.

\begin{env}[20.4.1]
\label{0.20.4.1-fr}
Soient $A$ un anneau topologique, $B$ une $A$-algèbre topologique ; la
$A$-algèbre $B\otimes_A B$ sera munie de la topologie \emph{produit
tensoriel}, qui en fait une $A$-algèbre topologique ; nous désignerons par $p$
(ou $p_{B/A}$) le $A$-homomorphisme canonique surjectif
\begin{equation}
\tag{20.4.1.1}
\label{0.20.4.1.1-fr}
  p:B\otimes_A B\to B
\end{equation}
tel que $p(b\otimes b')=bb'$ ; il est immédiat que $p$ est continu. Le
\emph{noyau} de $p$ sera noté $\mathfrak{J}_{B/A}$ (ou seulement
$\mathfrak{J}$ s'il n'y a pas de confusion). On désignera par
$j_1:B\to B\otimes_A B$ et $j_2:B\to B\otimes_A B$ les deux
$A$-homomorphismes canoniques, tels que
\[
  j_1(b)=b\otimes1, \qquad j_2(b)=1\otimes b
\]
qui sont continus.
\end{env}

\oldpage[0\textsubscript{IV-1}]{124}
\begin{definition}[20.4.2]
\label{0.20.4.2-fr}
On appelle \emph{$B$-algèbre augmentée des parties principales d'ordre $1$ de
$B$ par rapport à $A$} et l'on note $\mathrm{P}^1_{B/A}$ la $A$-algèbre
topologique quotient
\begin{equation}
\tag{20.4.2.1}
\label{0.20.4.2.1-fr}
  \mathrm{P}^1_{B/A}=(B\otimes_A B)/\mathfrak{J}^2
\end{equation}
munie de la structure de $B$-algèbre topologique définie par l'homomorphisme
$\bar{j}_1:B\to\mathrm{P}^1_{B/A}$ (déduit de $j_1$ par composition avec
l'homomorphisme canonique $B\otimes_A B\to\mathrm{P}^1_{B/A}$), et de
l'augmentation de $B$-algèbre $\varepsilon:\mathrm{P}^1_{B/A}\to B$ (aussi
notée $\varepsilon_{B/A}$) déduite de $p$ par passage au quotient.

Comme $p(b\otimes1)=b$ par définition, il est clair que $\varepsilon$ est
bien une augmentation de $B$-algèbre.
\end{definition}

\begin{definition}[20.4.3]
\label{0.20.4.3-fr}
Le noyau de l'augmentation $\varepsilon:\mathrm{P}^1_{B/A}\to B$,
\begin{equation}
\tag{20.4.3.1}
\label{0.20.4.3.1-fr}
  \Omega^1_{B/A}=\mathfrak{J}_{B/A}/(\mathfrak{J}_{B/A})^2
\end{equation}
muni de la topologie induite par celle de $\mathrm{P}^1_{B/A}$, qui en fait
un $B$-module topologique, est appelé le $B$-module des
\emph{$1$-différentielles} (ou simplement des \emph{différentielles}) de $B$
par rapport à $A$.
\end{definition}

On notera que la topologie de $\Omega^1_{B/A}$ est aussi la topologie quotient
de la topologie induite sur $\mathfrak{J}_{B/A}$ par celle de
$B\otimes_A B$ (Bourbaki, \emph{Top. gén.}, chap.~III, 3\textsuperscript{e}
éd., §~2, n\textsuperscript{o}~7, prop.~20). Si $B$ est discret il en est de
même de $\Omega^1_{B/A}$. On note $\widehat{\Omega}^1_{B/A}$ le
\emph{séparé complété} du $B$-module topologique $\Omega^1_{B/A}$.

Tout anneau topologique $B$ pouvant être considéré comme une
$\mathbf{Z}$-algèbre topologique (où $\mathbf{Z}$ est muni de la topologie
discrète), on peut définir le $B$-module topologique $\Omega^1_{B/\mathbf{Z}}$,
qu'on appelle parfois aussi le $B$-module des \emph{différentielles absolues}
sur $B$ et que l'on note $\Omega^1_B$. Si $B$ est une algèbre topologique sur
un corps premier $P$ (discret), on a
$B\otimes_P B=B\otimes_{\mathbf{Z}}B$, donc
$\Omega^1_{B/P}=\Omega^1_B$.

\begin{lemma}[20.4.4]
\label{0.20.4.4-fr}
Soient $A$ un anneau, $B$ une $A$-algèbre. L'idéal $\mathfrak{J}_{B/A}$ de
$B\otimes_A B$ est engendré par les éléments $1\otimes s-s\otimes1$, où $s$
parcourt un ensemble de générateurs de la $A$-algèbre $B$.
\end{lemma}

\begin{proof}
Il est clair que pour tout $x\in B$, on a
$x\otimes1-1\otimes x\in\mathfrak{J}$ ; d'autre part, quels que soient $x$,
$y$ dans $B$, on a
$x\otimes y=xy\otimes1+(x\otimes1)(1\otimes y-y\otimes1)$. Si
$\sum_i(x_i\otimes y_i)\in\mathfrak{J}$, on a par définition
$\sum_i x_i y_i=0$, donc
\begin{equation}
\tag{20.4.4.1}
\label{0.20.4.4.1-fr}
  \sum_i(x_i\otimes y_i)=
  \sum_i(x_i\otimes1)(1\otimes y_i-y_i\otimes1)
\end{equation}
ce qui prouve que $\mathfrak{J}$ est l'idéal engendré par les éléments
$1\otimes x-x\otimes1$. En outre, si $x=st$, on a
\begin{equation}
\tag{20.4.4.2}
\label{0.20.4.4.2-fr}
  x\otimes1-1\otimes x=(s\otimes1)(t\otimes1-1\otimes t)
  +(s\otimes1-1\otimes s)(1\otimes t)
\end{equation}
ce qui termine aussitôt la démonstration par récurrence.
\end{proof}

\begin{proposition}[20.4.5]
\label{0.20.4.5-fr}
Soient $A$ un anneau topologique, $B$ une $A$-algèbre topologique. La
topologie de $\Omega^1_{B/A}$ est moins fine que la topologie déduite de celle
de $B$ (19.0.2) ; si dans $B$ le carré de tout idéal ouvert est ouvert, ces
deux topologies sont identiques.
\end{proposition}

\begin{proof}
La première assertion est triviale, la topologie produit tensoriel sur
$B\otimes_A B$ étant moins fine que la topologie déduite de celle de $B$ ;
\emph{a fortiori} la topologie induite sur
$\mathfrak{J}=\mathfrak{J}_{B/A}$

\oldpage[0\textsubscript{IV-1}]{125}
par celle de $B\otimes_A B$ est moins fine que la topologie sur $\mathfrak{J}$
déduite de celle de $B$. Pour démontrer la seconde assertion, écrivons
$M\otimes N$, par abus de notation, le sous-module
$\operatorname{Im}(M\otimes_A N)$ pour deux sous-$A$-modules $M$, $N$ de $B$.
Utilisant la relation
\[
\begin{aligned}
  (xy)\otimes z-x\otimes(yz)
  &=(x\otimes1)(1\otimes z)(y\otimes1-1\otimes y)\\
  &=xz.(y\otimes1-1\otimes y)
    -x.(z\otimes1-1\otimes z)(y\otimes1-1\otimes y)
\end{aligned}
\]
dans le $B$-module $B\otimes_A B$ (défini par $j_1$), on voit aussitôt,
compte tenu de \sref{0.20.4.4-fr}, que, si $\mathfrak{R}$ est un idéal de
$B$, on a
\[
  ((\mathfrak{R}^2\otimes B)+(B\otimes\mathfrak{R}^2))\cap\mathfrak{J}
  \subset(\mathfrak{R}\otimes\mathfrak{R})\cap\mathfrak{J}
  +\mathfrak{R}\mathfrak{J}+\mathfrak{J}^2
\]
et d'autre part, si $x_i$, $y_i$ sont des éléments de $\mathfrak{R}$ tels
que $\sum_i(x_i\otimes y_i)\in\mathfrak{J}$, il résulte de
\sref{0.20.4.4.1-fr} que l'on a
$\sum_i(x_i\otimes y_i)\in\mathfrak{R}\mathfrak{J}$, si bien que finalement
\begin{equation}
\tag{20.4.5.1}
\label{0.20.4.5.1-fr}
  (\mathfrak{R}^2\otimes B+B\otimes\mathfrak{R}^2)\cap\mathfrak{J}
  \subset\mathfrak{R}\mathfrak{J}+\mathfrak{J}^2.
\end{equation}
Or, on a un système fondamental de voisinages de $0$ dans $\mathfrak{J}$
(pour la topologie induite par celle de $B\otimes_A B$) en prenant pour
voisinages de $0$ les ensembles
$(\mathfrak{R}\otimes B+B\otimes\mathfrak{R})\cap\mathfrak{J}$, où
$\mathfrak{R}$ parcourt l'ensemble des idéaux ouverts. Comme la topologie de
$\Omega^1_{B/A}$ déduite de celle de $B$ est aussi le quotient par
$\mathfrak{J}^2$ de la topologie de $\mathfrak{J}$ déduite de celle de $B$,
l'hypothèse sur les idéaux ouverts de $B$ et la relation
\sref{0.20.4.5.1-fr} achèvent de prouver la seconde assertion.
\end{proof}

\begin{definition}[20.4.6]
\label{0.20.4.6-fr}
Soient $p_1$ et $p_2$ les applications composées
$B\xrightarrow{j_1}B\otimes_A B\to\mathrm{P}^1_{B/A}$,
$B\xrightarrow{j_2}B\otimes_A B\to\mathrm{P}^1_{B/A}$, qui sont des
$A$-homomorphismes continus tels que
$\varepsilon\circ p_1=\varepsilon\circ p_2=\mathrm{I}_B$. Le
$A$-homomorphisme continu de $A$-modules
\begin{equation}
\tag{20.4.6.1}
\label{0.20.4.6.1-fr}
  d_{B/A}=p_1-p_2:B\to\Omega^1_{B/A}
\end{equation}
est appelé la \emph{différentielle extérieure de $B$ relative à $A$} ; pour
tout $x\in B$, $d_{B/A}(x)$ (aussi noté $d(x)$ ou $dx$) est appelé la
\emph{différentielle de $x$} (relativement à $A$).
\end{definition}

Lorsque $A=\mathbf{Z}$, on écrit $d_B$ au lieu de $d_{B/\mathbf{Z}}$ ; si $B$
est une algèbre sur un corps premier $P$, on a $d_B=d_{B/P}$.

\begin{proposition}[20.4.7]
\label{0.20.4.7-fr}
Le $B$-module $\Omega^1_{B/A}$ est engendré par les éléments $d_{B/A}(x)$,
où $x$ parcourt un système de générateurs de la $A$-algèbre $B$.
\end{proposition}

\begin{proof}
Comme $d_{B/A}(x)$ est l'image canonique de $x\otimes1-1\otimes x$ dans
$\Omega^1_{B/A}$, la proposition est une conséquence immédiate de
\sref{0.20.4.4-fr}.
\end{proof}

\begin{theorem}[20.4.8]
\label{0.20.4.8-fr}
Soient $A$ un anneau topologique, $B$ une $A$-algèbre topologique.
\begin{enumerate}
  \item[\rm(i)] Il existe un isomorphisme unique de $B$-algèbres augmentées
  topologiques
\begin{equation}
\tag{20.4.8.1}
\label{0.20.4.8.1-fr}
  \varphi:\mathrm{P}^1_{B/A}\xrightarrow{\sim}D_B(\Omega^1_{B/A})
\end{equation}
qui se réduit à l'identité dans $\Omega^1_{B/A}$.
  \item[\rm(ii)] L'homomorphisme $d_{B/A}$ est une $A$-dérivation de $B$
  dans $\Omega^1_{B/A}$, ayant la propriété universelle suivante : pour tout
  $B$-module topologique $L$, l'application $u\mapsto u\circ d_{B/A}$ est
  un isomorphisme de $A$-modules
\begin{equation}
\tag{20.4.8.2}
\label{0.20.4.8.2-fr}
  \operatorname{Hom.\ cont}_B(\Omega^1_{B/A},L)
  \xrightarrow{\sim}\operatorname{D\acute{e}r.\ cont}_A(B,L).
\end{equation}
\end{enumerate}
\end{theorem}

\begin{proof}
\oldpage[0\textsubscript{IV-1}]{126}
(i) Il est immédiat que $\varphi$ (avec les notations de
\sref{0.20.4.6-fr}) est nécessairement l'application
$z\mapsto(\varepsilon(z),z-p_1(\varepsilon(z)))$, et l'isomorphisme
réciproque l'application $(b,x)\mapsto p_1(b)+x$ ; ces deux applications
étant continues, cela prouve la première assertion. On notera que cela
entraîne que la topologie de $B$ s'identifie (par $p_1$) au \emph{quotient}
de la topologie de $B\otimes_A B$ par $\mathfrak{J}_{B/A}$.

(ii) Le fait que $d_{B/A}$ soit une $A$-dérivation de $B$ résulte de la
définition \sref{0.20.4.6-fr} et de \sref{0.20.1.1-fr}.

Pour prouver la propriété universelle, rappelons que
$\operatorname{D\acute{e}r.\ cont}_A(B,L)$ s'identifie canoniquement à
l'ensemble $G$ des homomorphismes continus de $A$-algèbres
$u:B\to D_B(L)$ tels que le composé
$B\xrightarrow{u}D_B(L)\xrightarrow{q}B$ soit l'identité (où $q(b,0)=b$ est
la projection ; cf. \sref{0.20.1.6-fr} et \sref{0.20.3.2-fr}). D'autre part,
grâce à l'isomorphisme $\varphi$,
$\operatorname{Hom.\ cont}_B(\Omega^1_{B/A},L)$ s'identifie canoniquement à
l'ensemble des homomorphismes continus de $B$-algèbres
$v:\mathrm{P}^1_{B/A}\to D_B(L)$ tels que le composé
$\mathrm{P}^1_{B/A}\xrightarrow{v}D_B(L)\xrightarrow{q}B$ soit
l'augmentation $\varepsilon$. Comme on a $p_2=p_1-d_{B/A}$ par définition,
tout revient à prouver que tout $u\in G$ se factorise en
\[
\xymatrix{
  B \ar[r]^u \ar[d]_{p_2} & D_B(L) \\
  \mathrm{P}^1_{B/A} \ar[ur]_v &
}
\]
où $v$ est un $B$-homomorphisme continu. Or, on a déjà un homomorphisme
continu de $A$-algèbres $j:b\mapsto(b,0)$ de $B$ dans $D_B(L)$, qui
appartient à $G$ ; par définition du produit tensoriel topologique
d'algèbres topologiques ($\mathbf{0_I}$, 7.7.6), il existe donc un
$A$-homomorphisme continu d'algèbres
$w:B\otimes_A B\to D_B(L)$ rendant commutatif le diagramme
\[
\xymatrix{
  B\otimes_A B \ar[dr]^w & B \ar[l]_{j_1} \ar[d]^j \\
  B \ar[u]^{j_2} \ar[r]_u & D_B(L)
}
\]
On a donc par définition
$w(b\otimes1-1\otimes b)=j(b)-u(b)\in L$, et en vertu de
\sref{0.20.4.4-fr}, cela entraîne $w(\mathfrak{J})\subset L$ donc
$w(\mathfrak{J}^2)=0$ ; par suite $w$ se factorise en
\[
  B\otimes_A B\to\mathrm{P}^1_{B/A}\xrightarrow{v}D_B(L)
\]
où $v$ est un homomorphisme continu de $A$-algèbres ; d'ailleurs, comme
$v\circ p_1=j$ est un homomorphisme de $B$-algèbres, il en est de même de
$v$ par définition de la structure de $B$-algèbre de
$\mathrm{P}^1_{B/A}$ ; comme on a par définition $u=v\circ p_2$, cela
termine la démonstration.
\end{proof}

\begin{theorem}[20.4.9]
\label{0.20.4.9-fr}
Supposons que $B$ soit une $A$-algèbre topologique formellement lisse. Alors
le $B$-module topologique $\Omega^1_{B/A}$ est \emph{formellement projectif}.
\end{theorem}

\begin{proof}
En effet, l'hypothèse entraîne que $B\otimes_A B$ (muni de la structure de
$B$-algèbre topologique définie par $j_1$) est une $B$-algèbre topologique
formellement lisse \sref{0.19.3.5-fr}, (iii), et par suite aussi une
$A$-algèbre topologique formellement lisse \sref{0.19.3.5-fr}, (ii) ; comme
$B$ est topologiquement isomorphe à la $A$-algèbre quotient
$(B\otimes_A B)/\mathfrak{J}$ et est une $A$-algèbre formellement lisse, la
conclusion résulte de \sref{0.19.5.3-fr}.
\end{proof}

\oldpage[0\textsubscript{IV-1}]{127}
\begin{corollary}[20.4.10]
\label{0.20.4.10-fr}
Supposons en outre que dans $B$ le carré de tout idéal ouvert soit ouvert ;
alors, pour tout idéal ouvert $\mathfrak{R}$ de $B$,
$\Omega^1_{B/A}\otimes_B(B/\mathfrak{R})
=\Omega^1_{B/A}/\mathfrak{R}.\Omega^1_{B/A}$ est un
$(B/\mathfrak{R})$-module projectif.
\end{corollary}

\begin{proof}
En effet, la topologie de $\Omega^1_{B/A}$ est alors déduite de celle de
$B$ \sref{0.20.4.5-fr}, et il suffit d'appliquer \sref{0.19.2.4-fr}.
\end{proof}

\begin{corollary}[20.4.11]
\label{0.20.4.11-fr}
Soient $A$, $B$ deux anneaux locaux noethériens, $\rho:A\to B$ un
homomorphisme local, faisant de $B$ une $A$-algèbre topologique formellement
lisse (pour les topologies préadiques). Alors, pour tout idéal de définition
$\mathfrak{b}$ de $B$,
$\Omega^1_{B/A}\otimes_B(B/\mathfrak{b})$ est un
$(B/\mathfrak{b})$-module libre.
\end{corollary}

\begin{proof}
En effet, il résulte de \sref{0.20.4.10-fr} que ce module est projectif, et
comme $B/\mathfrak{b}$ est un anneau artinien, tout
$(B/\mathfrak{b})$-module projectif est libre
($\mathbf{0_{III}}$, 10.1.3).
\end{proof}

\begin{proposition}[20.4.12]
\label{0.20.4.12-fr}
Soient $A$ un anneau topologique, $B$ une $A$-algèbre topologique. Si
l'homomorphisme structural $A\to B$ est surjectif, on a
$\Omega^1_{B/A}=0$.
\end{proposition}

\begin{proof}
En effet, on a $\operatorname{D\acute{e}r}_A(B,L)=0$ pour tout $B$-module
$L$ en vertu de \sref{0.20.1.1-fr}, et la proposition découle donc aussitôt
de \sref{0.20.4.8-fr}.
\end{proof}

\begin{examples}[20.4.13]
\label{0.20.4.13-fr}
\begin{enumerate}
  \item[\rm(i)] Soient $A$ un anneau,
  $B=A[X_\alpha]_{\alpha\in I}$ une algèbre de polynômes sur $A$. Alors
  $\Omega^1_{B/A}$ est un $B$-module libre, dont les $dX_\alpha$ forment une
  base. En effet, les $dX_\alpha$ engendrent ce $B$-module
  \sref{0.20.4.7-fr}. D'autre part, si $L$ est un $B$-module libre, ayant
  une base $(z_\alpha)_{\alpha\in I}$ dont $I$ est l'ensemble d'indices, il
  existe un $A$-homomorphisme $u$ de $B$ dans $D_B(L)$ tel que
  $u(X_\alpha)=(X_\alpha,z_\alpha)$, pour tout $\alpha\in I$, donc
  \sref{0.20.1.5-fr} une $A$-dérivation $D$ de $B$ dans $L$ telle que
  $D(X_\alpha)=z_\alpha$ pour tout $\alpha$ ; en vertu de
  \sref{0.20.4.8.1-fr}, cela prouve que les $dX_\alpha$ sont linéairement
  indépendants.
  \item[\rm(ii)] Soient $A$ un anneau, $L$ un $A$-module, $B$ la
  $A$-algèbre $D_A(L)$ ; alors l'homomorphisme canonique
  $x\mapsto d_{B/A}x$ de $L$ dans $\Omega^1_{B/A}$ est \emph{bijectif}, car
  il est immédiat que les $B$-dérivations de $B=D_A(L)$ dans un $B$-module
  $M$ sont les applications de la forme $(b,x)\mapsto u(x)$, où
  $u\in\operatorname{Hom}_A(L,M)$ ; on conclut donc par
  \sref{0.20.4.8-fr}, (ii).
  \item[\rm(iii)] Supposons que $B$ soit le produit de deux $A$-algèbres
  topologiques $B_1$, $B_2$ (identifiées à des idéaux de $B$). Alors les
  images de $B_1\otimes_A B_2$ et de $B_2\otimes_A B_1$ par l'homomorphisme
  $p$ \sref{0.20.4.1.1-fr} sont nulles, d'où il résulte aussitôt que
  $\mathrm{P}^1_{B/A}$ s'identifie au \emph{produit}
  $\mathrm{P}^1_{B_1/A}\times\mathrm{P}^1_{B_2/A}$, et que le $B$-module
  $\Omega^1_{B/A}$ est \emph{somme directe} (topologique) des $B$-modules
  $\Omega^1_{B_1/A}$ et $\Omega^1_{B_2/A}$ (annulés respectivement par
  $B_2$ et $B_1$).
  \item[\rm(iv)] Soient $A$, $B$ deux anneaux intègres tels que $A\subset
  B$, que $A$ soit intégralement clos, $B$ entier sur $A$, et que le corps
  des fractions de $B$ soit une extension séparable de celui de $A$. Alors
  le $B$-module $\Omega^1_{B/A}$ est un \emph{module de torsion}. En effet,
  pour tout $x\in B$, le polynôme minimal de $x$ par rapport au corps des
  fractions $K$ de $A$ est un polynôme $f(T)$ appartenant à $A[T]$ ; comme
  $x$ est séparable sur $K$, on a $f'(x)\neq0$ et d'autre part on déduit de
  la relation $f(x)=0$ que $f'(x)d_{B/A}x=0$, d'où notre assertion en vertu
  de \sref{0.20.4.7-fr}.
\end{enumerate}
\end{examples}

\begin{remarks}[20.4.14]
\label{0.20.4.14-fr}
\begin{enumerate}
  \item[\rm(i)] On notera que le $B$-module $\Omega^1_{B/A}$, privé de sa
  topologie, est indépendant des topologies de $A$ et de $B$.
  \item[\rm(ii)] Nous introduirons plus tard l'« algèbre des parties
  principales d'ordre $n$ » de $B$ par rapport à $A$,
  $\mathrm{P}^n_{B/A}=(B\otimes_A B)/(\mathfrak{J}_{B/A})^{n+1}$, qui est à
  la base du « calcul différentiel d'ordre $n$ ».
\end{enumerate}
\end{remarks}

\subsection{Propriétés fonctorielles fondamentales de
\texorpdfstring{$\Omega^1_{B/A}$}{Omega1(B/A)}}
\label{subsection:0.20.5-fr}

\oldpage[0\textsubscript{IV-1}]{128}
\begin{env}[20.5.1]
\label{0.20.5.1-fr}
Dans tout ce numéro et le suivant, sauf mention expresse du contraire, les
anneaux et modules considérés sont supposés munis de la topologie
\emph{discrète}.
\end{env}

\begin{env}[20.5.2]
\label{0.20.5.2-fr}
Soient $A$ un anneau, $B$, $C$ deux $A$-algèbres, $u:B\to C$ un
$A$-homomorphisme ; on a un diagramme commutatif
\[
\xymatrix{
  B\otimes_A B \ar[r]^{u\otimes u} \ar[d]_{p_{B/A}} &
  C\otimes_A C \ar[d]^{p_{C/A}} \\
  B \ar[r]^u & C
}
\tag{20.5.2.1}
\label{0.20.5.2.1-fr}
\]
d'où par passage aux quotients, un $A$-homomorphisme d'algèbres
\begin{equation}
\tag{20.5.2.2}
\label{0.20.5.2.2-fr}
  u':\mathrm{P}^1_{B/A}\to\mathrm{P}^1_{C/A}
\end{equation}
tel que le diagramme
\[
\xymatrix{
  \mathrm{P}^1_{B/A} \ar[r]^{u'} & \mathrm{P}^1_{C/A} \\
  B \ar[u]^{p_1} \ar[r]^u & C \ar[u]_{p_1}
}
\]
soit commutatif ; comme $u\otimes u$ applique $\mathfrak{J}_{B/A}$ dans
$\mathfrak{J}_{C/A}$, on obtient, en restreignant $u'$ à
$\Omega^1_{B/A}$, une application
\begin{equation}
\tag{20.5.2.3}
\label{0.20.5.2.3-fr}
  u'':\Omega^1_{B/A}\to\Omega^1_{C/A}
\end{equation}
telle que le couple $(u'',u)$ soit un \emph{di-homomorphisme} pour la
structure de $B$-module de $\Omega^1_{B/A}$ et la structure de $C$-module de
$\Omega^1_{C/A}$ ; ce dernier fait permet d'en déduire canoniquement un
\emph{homomorphisme de $C$-modules}
\begin{equation}
\tag{20.5.2.4}
\label{0.20.5.2.4-fr}
  u_{C/B/A}:\Omega^1_{B/A}\otimes_B C\to\Omega^1_{C/A}.
\end{equation}

\oldpage[0\textsubscript{IV-1}]{129}
En outre, comme le diagramme
\[
\xymatrix{
  \mathrm{P}^1_{B/A} \ar[r]^{u'} & \mathrm{P}^1_{C/A} \\
  B \ar[u]^{p_2} \ar[r]^u & C \ar[u]_{p_2}
}
\tag{20.5.2.5}
\label{0.20.5.2.5-fr}
\]
est aussi commutatif, on en déduit que le diagramme
\[
\xymatrix{
  \Omega^1_{B/A} \ar[r]^{u''} & \Omega^1_{C/A} \\
  B \ar[u]^{d_{B/A}} \ar[r]^u & C \ar[u]_{d_{C/A}}
}
\tag{20.5.2.6}
\label{0.20.5.2.6-fr}
\]
est commutatif.

Enfin, si $w:C\to D$ est un second homomorphisme de $A$-algèbres, on a la
propriété de transitivité
\begin{equation}
\tag{20.5.2.7}
\label{0.20.5.2.7-fr}
  (w\circ u)_{D/B/A}
  =w_{D/C/A}\circ(u_{C/B/A}\otimes\mathrm{I}_D)
\end{equation}
comme il résulte de la définition.
\end{env}

\begin{env}[20.5.3]
\label{0.20.5.3-fr}
Soient maintenant $A$, $B$ deux anneaux, $v:A\to B$ un homomorphisme
d'anneaux, $C$ une $B$-algèbre qui devient une $A$-algèbre au moyen de $v$ ;
alors l'application canonique $v_0:C\otimes_A C\to C\otimes_B C$ est un
\emph{di-homomorphisme} surjectif d'algèbres (relatif à $v:A\to B$) tel que
le diagramme
\[
\xymatrix{
  C\otimes_A C \ar[r]^{v_0} \ar[d]_{p_{C/A}} &
  C\otimes_B C \ar[d]^{p_{C/B}} \\
  C \ar[r]^{\mathrm{I}_C} & C
}
\tag{20.5.3.1}
\label{0.20.5.3.1-fr}
\]
soit commutatif ; par passage aux quotients, on en déduit un
di-homomorphisme d'algèbres
\begin{equation}
\tag{20.5.3.2}
\label{0.20.5.3.2-fr}
  v':\mathrm{P}^1_{C/A}\to\mathrm{P}^1_{C/B}
\end{equation}
tel que le diagramme
\[
\xymatrix{
  \mathrm{P}^1_{C/A} \ar[r]^{v'} & \mathrm{P}^1_{C/B} \\
  C \ar[u]^{p_1} \ar[r]^{\mathrm{I}_C} & C \ar[u]_{p_1}
}
\]
soit commutatif. Comme $v_0$ applique $\mathfrak{J}_{C/A}$ dans
$\mathfrak{J}_{C/B}$ on obtient, en restreignant $v'$ à
$\Omega^1_{C/A}$, une application
\begin{equation}
\tag{20.5.3.3}
\label{0.20.5.3.3-fr}
  v_{C/B/A}:\Omega^1_{C/A}\to\Omega^1_{C/B}
\end{equation}
qui est un homomorphisme de $C$-modules.

\oldpage[0\textsubscript{IV-1}]{130}
En outre, comme le diagramme
\[
\xymatrix{
  \mathrm{P}^1_{C/A} \ar[r]^{v'} & \mathrm{P}^1_{C/B} \\
  C \ar[u]^{p_2} \ar[r]^{\mathrm{I}_C} & C \ar[u]_{p_2}
}
\tag{20.5.3.4}
\label{0.20.5.3.4-fr}
\]
est aussi commutatif, on en déduit que le diagramme
\[
\xymatrix{
  \Omega^1_{C/A} \ar[r]^{v_{C/B/A}} & \Omega^1_{C/B} \\
  C \ar[u]^{d_{C/A}} \ar[r]^{\mathrm{I}_C} & C \ar[u]_{d_{C/B}}
}
\tag{20.5.3.5}
\label{0.20.5.3.5-fr}
\]
est commutatif.

Enfin, si $s:A'\to A$ est un second homomorphisme d'anneaux, on a la
propriété de transitivité
\begin{equation}
\tag{20.5.3.6}
\label{0.20.5.3.6-fr}
  (v\circ s)_{C/B/A'}=v_{C/B/A}\circ s_{C/A/A'}.
\end{equation}
\end{env}

\begin{env}[20.5.4]
\label{0.20.5.4-fr}
Si maintenant on a un diagramme commutatif d'homomorphismes d'anneaux
\[
\xymatrix{
  B \ar[r]^v & B' \\
  A \ar[u]^f \ar[r]^u & A' \ar[u]_g
}
\]
on déduit de \sref{0.20.5.2.4-fr} et \sref{0.20.5.3.3-fr}, par composition,
un homomorphisme de $B'$-modules
\begin{equation}
\tag{20.5.4.1}
\label{0.20.5.4.1-fr}
  \Omega^1_{B/A}\otimes_B B'
  \xrightarrow{v_{B'/B/A}}\Omega^1_{B'/A}
  \xrightarrow{u_{B'/A'/A}}\Omega^1_{B'/A'}
\end{equation}
tel que le diagramme de $A'$-homomorphismes
\[
\xymatrix{
  \Omega^1_{B/A}\otimes_B B' \ar[r] & \Omega^1_{B'/A'} \\
  B' \ar[u]^{d_{B/A}\otimes1} \ar[r]^{\mathrm{I}_{B'}} &
  B' \ar[u]_{d_{B'/A'}}
}
\tag{20.5.4.2}
\label{0.20.5.4.2-fr}
\]
soit commutatif.

L'homomorphisme \sref{0.20.5.4.1-fr} correspond d'ailleurs à un
di-homomorphisme de $B$-modules
\begin{equation}
\tag{20.5.4.3}
\label{0.20.5.4.3-fr}
  \Omega^1_{B/A}\to\Omega^1_{B'/A'}.
\end{equation}
\end{env}

\begin{proposition}[20.5.5]
\label{0.20.5.5-fr}
Si $A'$, $B$ sont deux $A$-algèbres et $B'=B\otimes_A A'$,
l'homomorphisme canonique \sref{0.20.5.4.1-fr}
\begin{equation}
\tag{20.5.5.1}
\label{0.20.5.5.1-fr}
  \Omega^1_{B/A}\otimes_B B'\to\Omega^1_{B'/A'}
\end{equation}
est bijectif.
\end{proposition}

\begin{proof}
\oldpage[0\textsubscript{IV-1}]{131}
Le premier membre de \sref{0.20.5.5.1-fr} n'est autre alors que
$\Omega^1_{B/A}\otimes_A A'$. On peut écrire
$B'\otimes_{A'}B'=(B\otimes_A B)\otimes_A A'$ à un isomorphisme canonique
près et $p_{B'/A'}$ s'identifie alors à
$p_{B/A}\otimes\mathrm{I}_{A'}$ ; par suite (comme $p_{B/A}$ est surjectif)
$\mathfrak{J}_{B'/A'}=\operatorname{Im}(\mathfrak{J}_{B/A}\otimes_A A')$
et
$\mathfrak{J}^2_{B'/A'}
=\operatorname{Im}(\mathfrak{J}^2_{B/A}\otimes_A A')$, d'où
$\mathrm{P}^1_{B'/A'}=\mathrm{P}^1_{B/A}\otimes_A A'$ à un isomorphisme
canonique près, qui transforme en eux-mêmes les idéaux d'augmentation ;
comme le $A$-module $\mathrm{P}^1_{B/A}$ s'identifie canoniquement à la
somme directe de $B$ et de $\Omega^1_{B/A}$, on a
\begin{equation}
\tag{20.5.5.2}
\label{0.20.5.5.2-fr}
  \Omega^1_{B/A}\otimes_A A'=\Omega^1_{B'/A'}
\end{equation}
par le même isomorphisme, et on vérifie aussitôt que le composé de cet
isomorphisme et de l'isomorphisme canonique
$\Omega^1_{B/A}\otimes_B B'\xrightarrow{\sim}
\Omega^1_{B/A}\otimes_A A'$ n'est autre que \sref{0.20.5.5.1-fr}.
\end{proof}

\begin{env}[20.5.6]
\label{0.20.5.6-fr}
Les homomorphismes canoniques \sref{0.20.5.2.4-fr} et
\sref{0.20.5.3.3-fr} donnent, par fonctorialité, pour tout $C$-module $L$,
des homomorphismes canoniques
\begin{align}
\tag{20.5.6.1}
\label{0.20.5.6.1-fr}
  &\operatorname{Hom}_C(\Omega^1_{C/A},L)
  \to\operatorname{Hom}_C(\Omega^1_{B/A}\otimes_B C,L)
  =\operatorname{Hom}_B(\Omega^1_{B/A},L) \\
\tag{20.5.6.2}
\label{0.20.5.6.2-fr}
  &\operatorname{Hom}_C(\Omega^1_{C/B},L)
  \to\operatorname{Hom}_C(\Omega^1_{C/A},L).
\end{align}
Compte tenu de \sref{0.20.4.8.2-fr} et des diagrammes commutatifs
\sref{0.20.5.2.6-fr} et \sref{0.20.5.3.5-fr}, ces homomorphismes ne sont
autres (à une identification canonique près) que les homomorphismes
\sref{0.20.2.1.2-fr} et \sref{0.20.2.1.3-fr} respectivement.
\end{env}

\begin{theorem}[20.5.7]
\label{0.20.5.7-fr}
Soient $u:A\to B$, $v:B\to C$ deux homomorphismes d'anneaux.
\begin{enumerate}
  \item[\rm(i)] La suite de $C$-modules
\begin{equation}
\tag{20.5.7.1}
\label{0.20.5.7.1-fr}
  \Omega^1_{B/A}\otimes_B C
  \xrightarrow{v_{C/B/A}}\Omega^1_{C/A}
  \xrightarrow{u_{C/B/A}}\Omega^1_{C/B}\to0
\end{equation}
est exacte.
  \item[\rm(ii)] Pour que $v_{C/B/A}$ soit inversible à gauche, il faut et
  il suffit que $C$ soit une $B$-algèbre formellement lisse relativement à
  $A$ (pour les topologies discrètes (cf. \sref{0.19.9.1-fr})) ; en
  particulier, il suffit pour cela que $C$ soit une $B$-algèbre formellement
  lisse (pour les topologies discrètes).
\end{enumerate}
\end{theorem}

\begin{proof}
(i) L'exactitude de la suite \sref{0.20.2.4.2-fr} montre tout d'abord,
compte tenu de \sref{0.20.5.6-fr}, que la suite
\[
  0\to\operatorname{Hom}_C(\Omega^1_{C/B},L)
  \to\operatorname{Hom}_C(\Omega^1_{C/A},L)
  \to\operatorname{Hom}_C(\Omega^1_{B/A}\otimes_B C,L)
\]
est exacte pour tout $C$-module $L$. On sait que cela implique l'exactitude
de la suite \sref{0.20.5.7.1-fr} (Bourbaki, \emph{Alg.}, chap.~II,
3\textsuperscript{e} éd., §~2, n\textsuperscript{o}~1, th.~1).

(ii) En vertu de l'exactitude de \sref{0.20.5.7.1-fr}, dire que
$v_{C/B/A}$ est inversible à gauche signifie que la suite
\begin{equation}
\tag{20.5.7.2}
\label{0.20.5.7.2-fr}
  0\to\Omega^1_{B/A}\otimes_B C
  \xrightarrow{v_{C/B/A}}\Omega^1_{C/A}
  \xrightarrow{u_{C/B/A}}\Omega^1_{C/B}\to0
\end{equation}
est \emph{exacte et scindée} ; on sait (Bourbaki, \emph{loc. cit.},
n\textsuperscript{o}~1, prop.~1) que cela équivaut à dire que pour tout
$C$-module $L$, la suite
\[
  0\to\operatorname{Hom}_C(\Omega^1_{C/B},L)
  \to\operatorname{Hom}_C(\Omega^1_{C/A},L)
  \to\operatorname{Hom}_C(\Omega^1_{B/A}\otimes_B C,L)\to0
\]
est exacte ; compte tenu de \sref{0.20.5.6-fr} et
\sref{0.20.2.4.2-fr}, cette condition équivaut à
\[
  \operatorname{Exalcom}_{B/A}(C,L)=0
\]
pour tout $C$-module $L$, et la conclusion résulte donc de
\sref{0.19.9.8.1-fr}.
\end{proof}

\oldpage[0\textsubscript{IV-1}]{132}
Notons en outre que si l'on a un diagramme commutatif d'homomorphismes
d'anneaux
\[
\xymatrix{
  A' \ar[r] & B' \ar[r] & C' \\
  A \ar[u] \ar[r] & B \ar[u] \ar[r] & C \ar[u]
}
\]
on en déduit un diagramme commutatif
\[
\xymatrix{
  \Omega^1_{B/A}\otimes_B C \ar[r] \ar[d] &
  \Omega^1_{C/A} \ar[r] \ar[d] &
  \Omega^1_{C/B} \ar[r] \ar[d] & 0 \\
  \Omega^1_{B'/A'}\otimes_{B'} C' \ar[r] &
  \Omega^1_{C'/A'} \ar[r] &
  \Omega^1_{C'/B'} \ar[r] & 0
}
\tag{20.5.7.3}
\label{0.20.5.7.3-fr}
\]
où les flèches verticales proviennent des di-homomorphismes
\sref{0.20.5.4.3-fr}.

\begin{corollary}[20.5.8]
\label{0.20.5.8-fr}
Supposons que l'homomorphisme $v:B\to C$ fasse de $C$ une $B$-algèbre
formellement étale (pour les topologies discrètes
\sref{0.19.10.2-fr}) ; alors l'homomorphisme \sref{0.20.5.3.3-fr}
\[
  v_{C/B/A}:\Omega^1_{B/A}\otimes_B C\to\Omega^1_{C/A}
\]
est bijectif.
\end{corollary}

\begin{proof}
En effet, si $C$ est une $B$-algèbre formellement non ramifiée pour les
topologies discrètes, il résulte de \sref{0.19.10.4-fr},
\sref{0.20.4.8-fr} et \sref{0.20.1.1-fr} que l'on a
$\operatorname{Hom}_C(\Omega^1_{C/B},L)=0$ pour tout $C$-module $L$, donc
$\Omega^1_{C/B}=0$ (cf. (20.7.4)) ; d'autre part, si $C$ est une
$B$-algèbre formellement lisse pour les topologies discrètes, la suite
\sref{0.20.5.7.2-fr} est exacte ; d'où le corollaire.
\end{proof}

\begin{corollary}[20.5.9]
\label{0.20.5.9-fr}
Soient $A$ un anneau, $B$ une $A$-algèbre, $S$ une partie multiplicative de
$B$ ; alors l'homomorphisme canonique
\begin{equation}
\tag{20.5.9.1}
\label{0.20.5.9.1-fr}
  S^{-1}\Omega^1_{B/A}\to\Omega^1_{S^{-1}B/A}
\end{equation}
est bijectif.
\end{corollary}

\begin{proof}
Il suffit d'appliquer \sref{0.20.5.8-fr} à $C=S^{-1}B$, qui est une
$B$-algèbre formellement étale pour les topologies discrètes
\sref{0.19.10.3-fr}, (ii).
\end{proof}

Compte tenu de \sref{0.20.5.5-fr}, on peut donc écrire
\begin{equation}
\tag{20.5.9.2}
\label{0.20.5.9.2-fr}
  \Omega^1_{S^{-1}B/S^{-1}A}
  =S^{-1}\Omega^1_{B/A}
  =\Omega^1_{S^{-1}B/A},
\end{equation}
à des isomorphismes canoniques près.

\begin{corollary}[20.5.10]
\label{0.20.5.10-fr}
Si $k$ est un corps et $K=k(X_\alpha)_{\alpha\in I}$ une extension pure de
$k$, les $dX_\alpha$ forment une base du $K$-espace vectoriel
$\Omega^1_{K/k}$.
\end{corollary}

\begin{proof}
Comme $K$ est le corps des fractions de l'anneau de polynômes
$k[X_\alpha]_{\alpha\in I}$, cela résulte de \sref{0.20.4.13-fr}, (i) et de
\sref{0.20.5.9-fr}.
\end{proof}

\oldpage[0\textsubscript{IV-1}]{133}
\begin{env}[20.5.11]
\label{0.20.5.11-fr}
Soient $A$ un anneau, $B$ une $A$-algèbre, $\mathfrak{R}$ un idéal de $B$,
$C$ la $A$-algèbre quotient $B/\mathfrak{R}$, et considérons
l'homomorphisme composé de $A$-modules
\begin{equation}
\tag{20.5.11.1}
\label{0.20.5.11.1-fr}
  \delta:\mathfrak{R}\longrightarrow B\xrightarrow{d_{B/A}}\Omega^1_{B/A}
\end{equation}
où la première flèche est l'injection canonique; comme
$d(xy)=xdy+ydx$, on voit que l'on a
$\delta(\mathfrak{R}^2)\subset\mathfrak{R}.\Omega^1_{B/A}$, d'où, par
passage aux quotients, un homomorphisme de $A$-modules
\begin{equation}
\tag{20.5.11.2}
\label{0.20.5.11.2-fr}
  \delta_{C/B/A}:\mathfrak{R}/\mathfrak{R}^2\longrightarrow
  \Omega^1_{B/A}\otimes_B C=
  \Omega^1_{B/A}/(\mathfrak{R}.\Omega^1_{B/A}).
\end{equation}
Mais en fait, $\delta_{C/B/A}$ est un homomorphisme de $C$-modules, car
pour $x\in B$ et $y\in\mathfrak{R}$, on a
$ydx\in\mathfrak{R}.\Omega^1_{B/A}$, donc
$d(xy)\equiv xdy\pmod{\mathfrak{R}.\Omega^1_{B/A}}$, ce qui prouve
d'abord que (20.5.11.2) est un homomorphisme de $B$-modules, et comme
$\mathfrak{R}$ annule les deux membres, cela établit notre assertion.

Si $B'$ est une seconde $A$-algèbre, $u:B\to B'$ un $A$-homomorphisme,
$\mathfrak{R}'$ un idéal de $B'$ tel que
$u(\mathfrak{R})\subset\mathfrak{R}'$ et $C'=B'/\mathfrak{R}'$ l'algèbre
quotient, on a un diagramme commutatif
\begin{equation}
\tag{20.5.11.3}
\label{0.20.5.11.3-fr}
\xymatrix{
  \mathfrak{R}/\mathfrak{R}^2 \ar[r]^-{\delta_{C/B/A}} \ar[d] &
  \Omega^1_{B/A}\otimes_B C \ar[d] \\
  \mathfrak{R}'/\mathfrak{R}'^2 \ar[r]_-{\delta_{C'/B'/A}} &
  \Omega^1_{B'/A}\otimes_{B'} C'
}
\end{equation}
où les flèches verticales proviennent de $u$ (20.5.2.4).
\end{env}

\begin{theorem}[20.5.12]
\label{0.20.5.12-fr}
Soient $A$ un anneau, $B$ une $A$-algèbre, $C$ la $A$-algèbre quotient
$B/\mathfrak{R}$, $u:B\to C$ l'homomorphisme canonique.
\begin{enumerate}
  \item[\rm(i)] On a une suite exacte de $C$-modules
\begin{equation}
\tag{20.5.12.1}
\label{0.20.5.12.1-fr}
  \mathfrak{R}/\mathfrak{R}^2
  \xrightarrow{\delta_{C/B/A}}\Omega^1_{B/A}\otimes_B C
  \xrightarrow{u_{C/B/A}}\Omega^1_{C/A}\longrightarrow 0,
\end{equation}
où $u_{C/B/A}$ et $\delta_{C/B/A}$ sont définis par (20.5.2.4) et
(20.5.11.2) respectivement.
  \item[\rm(ii)] Si l'on pose $E=B/\mathfrak{R}^2$, l'homomorphisme
canonique (20.5.2.4)
\[
  \Omega^1_{B/A}\otimes_B C\longrightarrow
  \Omega^1_{E/A}\otimes_E C
\]
est bijectif.
  \item[\rm(iii)] Les trois conditions suivantes sont équivalentes :
  \begin{enumerate}
    \item[\rm a)] $\delta_{C/B/A}$ est inversible à gauche.
    \item[\rm b)] Toute $A$-extension de $C$ par un $C$-module $L$, dont
l'image réciproque par $u:B\to C$ est $A$-triviale, est elle-même
$A$-triviale.
    \item[\rm c)] La $A$-algèbre $E=B/\mathfrak{R}^2$ est une extension
$A$-triviale de $C$ par $\mathfrak{R}/\mathfrak{R}^2$.
  \end{enumerate}
  \item[\rm(iv)] Il y a correspondance biunivoque canonique entre les
inverses à gauche de $\delta_{C/B/A}$ et les inverses à droite de
l'homomorphisme canonique $E\to C$.
\end{enumerate}
\end{theorem}

\begin{proof}
\oldpage[0\textsubscript{IV-1}]{134}
(i) Comme $u$ est surjectif, on a $\operatorname{D\acute{e}r}_B(C,L)=0$ pour tout $C$-module
$L$ en vertu de (20.1.1). La suite exacte (20.2.3.1) devient donc
\begin{equation}
\tag{20.5.12.2}
\label{0.20.5.12.2-fr}
  0\longrightarrow\operatorname{D\acute{e}r}_A(C,L)\xrightarrow{v^0}\operatorname{D\acute{e}r}_A(B,L)
  \xrightarrow{\partial}\operatorname{Exalcom}_B(C,L)
  \xrightarrow{v^1}\operatorname{Exalcom}_A(C,L)
  \xrightarrow{u^1}\operatorname{Exalcom}_A(B,L),
\end{equation}
où $v$ est l'homomorphisme $A\to B$. Rappelons d'autre part (18.3.8) que
$\operatorname{Exalcom}_B(C,L)$ s'identifie canoniquement à
$\operatorname{Hom}_C(\mathfrak{R}/\mathfrak{R}^2,L)$; on déduit donc de (20.5.12.2)
et (20.4.8) la suite exacte
\begin{equation}
\tag{20.5.12.3}
\label{0.20.5.12.3-fr}
\begin{split}
  0\longrightarrow\operatorname{Hom}_C(\Omega^1_{C/A},L)
  \longrightarrow\operatorname{Hom}_C(\Omega^1_{B/A}\otimes_B C,L)
  \xrightarrow{\varphi}\operatorname{Hom}_C(\mathfrak{R}/\mathfrak{R}^2,L)\\
  \xrightarrow{\psi}\operatorname{Ker}(u^1)\longrightarrow 0
\end{split}
\end{equation}
avec $\varphi=\partial\circ\eta^{-1}$ et $\psi=v^1\circ\eta$. Revenant
aux définitions de $\partial$ (20.2.2) et de $\eta$ (18.3.8), on constate
aussitôt que $\varphi$ est précisément l'homomorphisme
$\operatorname{Hom}(\delta_{C/B/A},1_L)$. L'existence de la suite exacte formée des
4 premiers termes de (20.5.12.3) montre donc que la suite (20.5.12.1)
est exacte (Bourbaki, \emph{Alg.}, chap. II, 3\textsuperscript{e} éd.,
\S~2, n\textsuperscript{o}~1, th.~1).

\medskip
(ii) Appliquons à $B$ et à l'idéal $\mathfrak{R}^2$ la suite exacte
(20.5.12.1), ce qui donne
\begin{equation}
\tag{20.5.12.4}
\label{0.20.5.12.4-fr}
  \mathfrak{R}^2/\mathfrak{R}^4\longrightarrow
  \Omega^1_{B/A}\otimes_B E\longrightarrow\Omega^1_{E/A}\longrightarrow 0,
\end{equation}
d'où, en tensorisant par $C$ (considéré comme $E$-algèbre), la suite
exacte
\[
  \mathfrak{R}^2/\mathfrak{R}^4\otimes_E C\longrightarrow
  \Omega^1_{B/A}\otimes_B C\longrightarrow
  \Omega^1_{E/A}\otimes_E C\longrightarrow 0.
\]
Or, si $x,y$ sont deux éléments de $\mathfrak{R}$, et $\zeta$ la classe
de $xy$ mod. $\mathfrak{R}^4$, l'image
$\delta'(\zeta\otimes1)$ est par définition
$d_{B/A}(xy)\otimes1=(xd_{B/A}(y)+yd_{B/A}(x))\otimes1$; mais comme les
images de $x$ et de $y$ dans $C$ sont nulles, on a aussi
$d_{B/A}(xy)\otimes1=0$ dans $\Omega^1_{B/A}\otimes_B C$, ce qui prouve
notre assertion.

\medskip
(iii) Dire que $\delta_{C/B/A}$ est inversible à gauche revient, compte
tenu de l'exactitude de (20.5.12.1), à dire que la suite
\begin{equation}
\tag{20.5.12.5}
\label{0.20.5.12.5-fr}
  0\longrightarrow\mathfrak{R}/\mathfrak{R}^2
  \xrightarrow{\delta_{C/B/A}}\Omega^1_{B/A}\otimes_B C
  \xrightarrow{u_{C/B/A}}\Omega^1_{C/A}\longrightarrow0
\end{equation}
est \emph{exacte et scindée}, et il revient au même (Bourbaki,
\emph{loc. cit.}) de dire que $\operatorname{Ker}(u^1)=0$ dans la suite exacte
(20.5.12.3) pour tout $L$, ce qui montre l'équivalence des conditions
a) et b) (cf. (18.3.6.2)).

Le fait que b) entraîne c) provient de ce que l'image réciproque par
$u:B\to C$ de la $A$-extension $E$ de $C$ par
$\mathfrak{R}/\mathfrak{R}^2$ est $A$-triviale, $u$ étant composée des
homomorphismes canoniques $B\to E\to C$ (18.1.6). Inversement, c)
entraîne b), car toute $B$-extension de $C$ par $L$ est $B$-équivalente
à $E\oplus_{\mathfrak{R}/\mathfrak{R}^2}L$ (18.3.8), autrement dit sa
classe est l'image de la classe de $E$ (considéré comme $B$-extension)
par l'homomorphisme

\oldpage[0\textsubscript{IV-1}]{135}
$w_*:\operatorname{Exan}_B(C,\mathfrak{R}/\mathfrak{R}^2)
\to\operatorname{Exan}_B(C,L)$ correspondant à un $B$-homomorphisme
$w:\mathfrak{R}/\mathfrak{R}^2\to L$. Le fait que c) entraîne b) résulte
alors de la commutativité du diagramme (18.3.6.5)
\[
\xymatrix{
  \operatorname{Exan}_B(C,\mathfrak{R}/\mathfrak{R}^2)
    \ar[r]^{w_*} \ar[d]_{u^*} &
  \operatorname{Exan}_B(C,L) \ar[d]^{u^*} \\
  \operatorname{Exan}_A(C,\mathfrak{R}/\mathfrak{R}^2)
    \ar[r]_{w_*} & \operatorname{Exan}_A(C,L).
}
\]

\medskip
(iv) On a vu (20.1.7) que les inverses à droite de $E\to C$
correspondent biunivoquement de façon canonique à l'ensemble des éléments
$D\in\operatorname{D\acute{e}r}_A(E,\mathfrak{R}/\mathfrak{R}^2)$ tels que $D(x)=x$ dans
$\mathfrak{R}/\mathfrak{R}^2$, donc aussi, par (20.4.8), à l'ensemble des
$E$-homomorphismes
$h:\Omega^1_{E/A}\to\mathfrak{R}/\mathfrak{R}^2$ tels que le composé
\[
  \mathfrak{R}/\mathfrak{R}^2\xrightarrow{d_{E/A}}
  \Omega^1_{E/A}\xrightarrow{h}\mathfrak{R}/\mathfrak{R}^2
\]
soit l'identité. Par tensorisation avec $C$, on en déduit (puisque
$\mathfrak{R}/\mathfrak{R}^2$ est un $C$-module) que le composé
\[
  \mathfrak{R}/\mathfrak{R}^2\xrightarrow{\delta_{C/E/A}}
  \Omega^1_{E/A}\otimes_E C\xrightarrow{h\otimes1}
  \mathfrak{R}/\mathfrak{R}^2
\]
est l'identité; or, puisque $\mathfrak{R}/\mathfrak{R}^2$ est un
$C$-module, $h\mapsto h\otimes1$ est un isomorphisme de l'ensemble
$\operatorname{Hom}_E(\Omega^1_{E/A},\mathfrak{R}/\mathfrak{R}^2)$ sur
$\operatorname{Hom}_C(\Omega^1_{E/A}\otimes_E C,\mathfrak{R}/\mathfrak{R}^2)$; et par
ailleurs (ii) prouve que l'on peut identifier canoniquement
$\Omega^1_{E/A}\otimes_E C$ et $\Omega^1_{B/A}\otimes_B C$,
$\delta_{C/E/A}$ s'identifiant alors à $\delta_{C/B/A}$. C.Q.F.D.
\end{proof}

\begin{example}[20.5.13]
\label{0.20.5.13-fr}
Soient $B=A[X_\alpha]_{\alpha\in I}$ une algèbre de polynômes sur $A$,
$\mathfrak{R}$ un idéal de $B$, $(P_\lambda)$ un système de générateurs
de $\mathfrak{R}$ et $C=B/\mathfrak{R}$; on sait que
$\Omega^1_{B/A}$ est un $B$-module libre dont les $dX_\alpha$ forment une
base (20.4.13, (i)), donc les $dX_\alpha$ forment aussi une base du
$C$-module libre $\Omega^1_{B/A}\otimes_B C$. D'autre part, il résulte
aussitôt de la définition que l'image de
$\mathfrak{R}/\mathfrak{R}^2$ par $\delta_{C/B/A}$ est le sous-$C$-module
engendré par les
\[
  dP_\lambda=\sum_\alpha
  \frac{\partial P_\lambda}{\partial X_\alpha}\,dX_\alpha.
\]
On en conclut que $\Omega^1_{C/A}$ est isomorphe au quotient du
$C$-module libre ayant les $dX_\alpha$ pour base, par le sous-$C$-module
engendré par les $dP_\lambda$, ce qui donne une description d'un module
de différentielles d'une algèbre \emph{quelconque}, toute $A$-algèbre
$C$ pouvant s'obtenir de la façon précédente.
\end{example}

\begin{corollary}[20.5.14]
\label{0.20.5.14-fr}
Si $C$ est une $A$-algèbre formellement lisse (pour les topologies
discrètes), la suite
\begin{equation}
\tag{20.5.14.1}
\label{0.20.5.14.1-fr}
  0\longrightarrow\mathfrak{R}/\mathfrak{R}^2
  \xrightarrow{\delta_{C/B/A}}\Omega^1_{B/A}\otimes_B C
  \xrightarrow{u_{C/B/A}}\Omega^1_{C/A}\longrightarrow0
\end{equation}
est exacte et scindée.
\end{corollary}

\begin{proof}
En effet, toute $A$-extension de $C$ par un $C$-module est alors triviale
(19.4.4.1).
\end{proof}

\oldpage[0\textsubscript{IV-1}]{136}
\begin{remark}[20.5.15]
\label{0.20.5.15-fr}
Soit $u:A\to B$ un homomorphisme \emph{surjectif} d'anneaux; alors, pour
tout homomorphisme d'anneaux $v:B\to C$, l'homomorphisme canonique
\begin{equation}
\tag{20.5.15.1}
\label{0.20.5.15.1-fr}
  u_{C/B/A}:\Omega^1_{C/A}\longrightarrow\Omega^1_{C/B}
\end{equation}
est \emph{bijectif}; cela résulte en effet de la suite exacte (20.5.7.1),
puisque $\Omega^1_{B/A}=0$ (20.4.12).
\end{remark}

\subsection{Modules d'imperfection et homomorphismes caractéristiques}
\label{subsection:0.20.6-fr}
\label{0.20.6-fr}

\begin{definition}[20.6.1]
\label{0.20.6.1-fr}
Étant donnés deux homomorphismes d'anneaux $u:A\to B$, $v:B\to C$, on
appelle \emph{module d'imperfection de la $B$-algèbre $C$ relativement à
$A$}, et on note $\Upsilon_{C/B/A}$, le $C$-module noyau de
l'homomorphisme
$v_{C/B/A}:\Omega^1_{B/A}\otimes_B C\to\Omega^1_{C/A}$.
On a donc par définition (cf. (20.5.7)) la suite exacte
\begin{equation}
\tag{20.6.1.1}
\label{0.20.6.1.1-fr}
  0\longrightarrow\Upsilon_{C/B/A}\longrightarrow
  \Omega^1_{B/A}\otimes_B C\xrightarrow{v_{C/B/A}}
  \Omega^1_{C/A}\xrightarrow{u_{C/B/A}}\Omega^1_{C/B}
  \longrightarrow0.
\end{equation}
\end{definition}

Lorsque $A=\mathbb{Z}$ (les modules $\Omega^1_{B/A}$ et $\Omega^1_{C/A}$
étant donc les modules de différentielles « absolues » $\Omega^1_B$ et
$\Omega^1_C$), on écrit $\Upsilon_{C/B}$ au lieu de
$\Upsilon_{C/B/\mathbb{Z}}$; lorsque $B$ et $C$ sont des algèbres sur un
corps premier $P$, on a $\Upsilon_{C/B/P}=\Upsilon_{C/B}$.

Soient $R$, $S$ des parties multiplicatives de $B$ et $C$ respectivement,
telles que l'image de $R$ soit contenue dans $S$. Il résulte alors de la
suite exacte (20.6.1.1) et de (20.5.9) que l'on a
\begin{equation}
\tag{20.6.1.2}
\label{0.20.6.1.2-fr}
  \Upsilon_{S^{-1}C/R^{-1}B/A}=S^{-1}\Upsilon_{C/B/A}.
\end{equation}

\begin{proposition}[20.6.2]
\label{0.20.6.2-fr}
Si $C$ est une $B$-algèbre formellement lisse relativement à $A$ (et en
particulier si $C$ est une $B$-algèbre formellement lisse), on a
$\Upsilon_{C/B/A}=0$.
\end{proposition}

\begin{proof}
Cela résulte de (20.5.7, (ii)).
\end{proof}

\begin{proposition}[20.6.3]
\label{0.20.6.3-fr}
Soient $k$ un corps, $K$ une extension de $k$. Pour que $K$ soit une
extension séparable de $k$, il faut et il suffit que $\Upsilon_{K/k}=0$
(autrement dit, l'homomorphisme canonique
$\Omega^1_k\otimes_k K\to\Omega^1_K$ est injectif ou encore (20.4.8),
toute dérivation de $k$ dans $K$ se prolonge en une dérivation de $K$
dans lui-même).
\end{proposition}

\begin{proof}
En effet, il est équivalent de dire que $K$ est séparable sur $k$ ou une
$k$-algèbre formellement lisse (19.6.1); d'autre part, si $P$ est le
corps premier de $k$, $K$ est séparable sur $P$, donc une $P$-algèbre
formellement lisse, et il revient donc au même de dire que $K$ est une
$k$-algèbre formellement lisse ou une $k$-algèbre formellement lisse
relativement à $P$. Enfin, dire que $K$ est une $k$-algèbre formellement
lisse relativement à $P$ équivaut d'après (20.5.7, (ii)) à dire que
l'homomorphisme
$v_{K/k/P}:\Omega^1_k\otimes_k K\to\Omega^1_K$ est inversible à gauche;
mais comme $K$ est un corps, cette dernière condition équivaut à dire
que le noyau de $v_{K/k/P}$, c'est-à-dire $\Upsilon_{K/k}$, est nul.
\end{proof}

\begin{env}[20.6.4]
\label{0.20.6.4-fr}
Considérons un diagramme commutatif
\begin{equation}
\tag{20.6.4.1}
\label{0.20.6.4.1-fr}
\xymatrix{
  A' \ar[r]^{u'} & B' \ar[r]^{v'} & C' \\
  A \ar[u]^f \ar[r]_u & B \ar[u]^g \ar[r]_v & C \ar[u]^h
}
\end{equation}

\oldpage[0\textsubscript{IV-1}]{137}
d'homomorphismes d'anneaux commutatifs. La commutativité du diagramme
correspondant (20.5.7.3) entraîne l'existence d'un unique
$C$-homomorphisme
\begin{equation}
\tag{20.6.4.2}
\label{0.20.6.4.2-fr}
  \Upsilon_{C/B/A}\longrightarrow\Upsilon_{C'/B'/A'}
\end{equation}
déduit canoniquement de (20.6.4.1) et rendant commutatif le diagramme
\begin{equation}
\tag{20.6.4.3}
\label{0.20.6.4.3-fr}
\xymatrix{
  0 \ar[r] & \Upsilon_{C/B/A} \ar[r] \ar[d] &
  \Omega^1_{B/A}\otimes_B C \ar[r]^{v_{C/B/A}} \ar[d] &
  \Omega^1_{C/A} \ar[r]^{u_{C/B/A}} \ar[d] &
  \Omega^1_{C/B} \ar[r] \ar[d] & 0 \\
  0 \ar[r] & \Upsilon_{C'/B'/A'} \ar[r] &
  \Omega^1_{B'/A'}\otimes_{B'} C' \ar[r]_{v_{C'/B'/A'}} &
  \Omega^1_{C'/A'} \ar[r]_{u_{C'/B'/A'}} &
  \Omega^1_{C'/B'} \ar[r] & 0
}
\end{equation}

La donnée de l'homomorphisme (20.6.4.2) équivaut d'ailleurs à celle d'un
$C'$-homomorphisme
\begin{equation}
\tag{20.6.4.4}
\label{0.20.6.4.4-fr}
  \Upsilon_{C/B/A}\otimes_C C'\longrightarrow\Upsilon_{C'/B'/A'}
\end{equation}
qui, composé avec l'homomorphisme canonique
$\Upsilon_{C/B/A}\to\Upsilon_{C/B/A}\otimes_C C'$, redonne (20.6.4.2).
Il est clair que (20.6.4.2) jouit d'une propriété évidente de
transitivité, permettant de dire que $\Upsilon_{C/B/A}$ est un
\emph{foncteur} en le triplet $(A,B,C)$.
\end{env}

\begin{env}[20.6.5]
\label{0.20.6.5-fr}
Il nous sera commode pour la suite, sous les conditions de (20.6.1),
d'introduire un \emph{complexe} (de chaînes) \emph{de $C$-modules}
$\mathrm{K}_\bullet(C/B/A)$ dont les termes sont nuls sauf en degrés 0
et 1, où on prend
\begin{equation}
\tag{20.6.5.1}
\label{0.20.6.5.1-fr}
\left\{
\begin{aligned}
  \mathrm{K}_0(C/B/A)&=\Omega^1_{C/A},\\
  \mathrm{K}_1(C/B/A)&=\Omega^1_{B/A}\otimes_B C,
\end{aligned}
\right.
\end{equation}
l'opérateur différentiel $\mathrm{K}_1\to\mathrm{K}_0$ étant
$v_{C/B/A}$. Cela permet d'écrire (à des isomorphismes canoniques près)
$\Omega^1_{C/B}$ et $\Upsilon_{C/B/A}$ comme les \emph{modules
d'homologie} de ce complexe
\begin{equation}
\tag{20.6.5.2}
\label{0.20.6.5.2-fr}
  \mathrm{H}_0(\mathrm{K}_\bullet(C/B/A))=\Omega^1_{C/B},\qquad
  \mathrm{H}_1(\mathrm{K}_\bullet(C/B/A))=\Upsilon_{C/B/A}.
\end{equation}
\end{env}

De même :
\begin{proposition}[20.6.6]
\label{0.20.6.6-fr}
Sous les hypothèses de (20.6.1), pour tout $C$-module $L$, on a des
$C$-isomorphismes canoniques
\begin{align}
\tag{20.6.6.1}
\label{0.20.6.6.1-fr}
  \mathrm{H}^0(\mathrm{K}_\bullet(C/B/A),L)&\simeq\operatorname{D\acute{e}r}_B(C,L),\\
\tag{20.6.6.2}
\label{0.20.6.6.2-fr}
  \mathrm{H}^1(\mathrm{K}_\bullet(C/B/A),L)&\simeq
  \operatorname{Exalcom}_{B/A}(C,L).
\end{align}
\end{proposition}

\begin{proof}
En effet, le complexe de cochaînes
$\operatorname{Hom}_C(\mathrm{K}_\bullet(C/B/A),L)$ n'est autre, en vertu de (20.4.8)
et (20.5.6), que le complexe
\[
  \cdots\longrightarrow0\longrightarrow\operatorname{D\acute{e}r}_A(C,L)
  \longrightarrow\operatorname{D\acute{e}r}_A(B,L)\longrightarrow0\longrightarrow\cdots
\]
\oldpage[0\textsubscript{IV-1}]{138}
où l'opérateur différentiel est $v^0$ (avec les notations de (20.2.1)).
La proposition résulte alors de la suite exacte (20.2.4.2) et de la
définition des modules de cohomologie
\begin{equation}
\tag{20.6.6.3}
\label{0.20.6.6.3-fr}
  \mathrm{H}^i(\mathrm{K}_\bullet,L)=
  \mathrm{H}^i(\operatorname{Hom}_C(\mathrm{K}_\bullet,L)).
\end{equation}
\end{proof}

Si on a un diagramme commutatif d'homomorphismes d'anneaux
\[
\xymatrix{
  A' \ar[r] & B' \ar[r] & C' \\
  A \ar[u] \ar[r] & B \ar[u] \ar[r] & C \ar[u]
}
\]
les di-homomorphismes (20.5.4.3) définissent un di-homomorphisme de
complexes de modules
\begin{equation}
\tag{20.6.6.4}
\label{0.20.6.6.4-fr}
  \mathrm{K}_\bullet(C/B/A)\longrightarrow
  \mathrm{K}_\bullet(C'/B'/A')
\end{equation}
et les di-homomorphismes qu'on en déduit pour l'homologie ou la
cohomologie s'identifient, par les formules (20.6.5.2), (20.6.6.1) et
(20.6.6.2), aux di-homomorphismes déjà définis dans (20.5.4.3),
(20.6.4.2), (20.2.1) et (18.3.7.2) (compte tenu, pour le dernier, de la
définition de l'opérateur $\partial$ dans la suite exacte (20.2.4.2)).

\begin{env}[20.6.7]
\label{0.20.6.7-fr}
On sait que pour un complexe $\mathrm{K}_\bullet$ de $C$-modules et un
$C$-module $L$, on a des homomorphismes canoniques
\[
  \alpha_i:\mathrm{H}^i(\mathrm{K}_\bullet,L)\longrightarrow
  \operatorname{Hom}_C(\mathrm{H}_i(\mathrm{K}_\bullet),L)
\]
(M, IV, 6). Ici, l'homomorphisme canonique
\[
  \alpha_1:\mathrm{H}^1(\mathrm{K}_\bullet(C/B/A),L)
  \longrightarrow
  \operatorname{Hom}_C(\mathrm{H}_1(\mathrm{K}_\bullet(C/B/A)),L)
\]
se définit immédiatement comme obtenu par passage au quotient par l'image
de $\operatorname{Hom}_C(\mathrm{K}_0,L)$ de l'homomorphisme de restriction
\[
  \operatorname{Hom}_C(\mathrm{K}_1(C/B/A),L)\longrightarrow
  \operatorname{Hom}_C(\mathrm{H}_1(\mathrm{K}_\bullet(C/B/A)),L)
\]
puisque $\mathrm{H}_1(\mathrm{K}_\bullet(C/B/A))$ n'est autre que le
noyau de $\mathrm{K}_1\to\mathrm{K}_0$; tenant compte de (20.6.6.2) et
(20.6.5.2), on obtient donc un $C$-homomorphisme canonique
\begin{equation}
\tag{20.6.7.1}
\label{0.20.6.7.1-fr}
  \operatorname{Exalcom}_{B/A}(C,L)\longrightarrow
  \operatorname{Hom}_C(\Upsilon_{C/B/A},L)
\end{equation}
qui s'explicite de la façon suivante : en vertu de (20.2.4.2), toute
$B$-extension de $C$ par $L$ qui est $A$-triviale provient de la donnée
d'une $A$-dérivation $D$ de $B$ dans $L$, donc (20.4.8) d'un
$C$-homomorphisme $f$ de $\Omega^1_{B/A}\otimes_B C$ dans $L$; on fait
correspondre à la classe de cette extension la \emph{restriction} de $f$
à $\Upsilon_{C/B/A}$, qui ne dépend effectivement que de cette classe et
non du choix de $D$.
\end{env}

\begin{definition}[20.6.8]
\label{0.20.6.8-fr}
Soient $u:A\to B$, $v:B\to C$ deux homomorphismes d'anneaux. Pour toute
$B$-extension $E$ de $C$ par un $C$-module $L$, qui est $A$-triviale
(18.3.7), on appelle \emph{homomorphisme caractéristique de $E$}, et on
note $\chi_E$, le $C$-homomorphisme $\Upsilon_{C/B/A}\to L$, image de la
classe de $E$ par l'homomorphisme canonique (20.6.7.1).
\end{definition}

\oldpage[0\textsubscript{IV-1}]{139}
On peut définir l'homomorphisme $\chi_E$ d'une autre manière :
\begin{proposition}[20.6.9]
\label{0.20.6.9-fr}
Soit $E$ une $B$-extension de $C$ par un $C$-module $L$, qui est
$A$-triviale (18.3.7); alors le diagramme
\begin{equation}
\tag{20.6.9.1}
\label{0.20.6.9.1-fr}
\xymatrix{
  0 \ar[r] & \Upsilon_{C/B/A} \ar[r] \ar[d]^{\chi_E} &
  \Omega^1_{B/A}\otimes_B C \ar[r]^{v_{C/B/A}}
    \ar[d]^{q_{E/B/A}\otimes1_C} &
  \Omega^1_{C/A} \ar[d]^{1} \\
  0 \ar[r] & L \ar[r]_{\delta_{C/E/A}} &
  \Omega^1_{E/A}\otimes_E C \ar[r]_{p_{C/E/A}} &
  \Omega^1_{C/A} \ar[r] & 0
}
\end{equation}
où $q:B\to E$ définit la structure de $B$-extension de $E$ et
$p:E\to C$ est l'homomorphisme d'augmentation, est commutatif et ses
lignes sont exactes.
\end{proposition}

\begin{proof}
La ligne inférieure du diagramme est la suite (20.5.12.5) relative aux
deux homomorphismes $A\xrightarrow{q\circ u}E$ et
$p:E\to C=E/L$; comme $p$ est surjectif et que $E$ est une extension
$A$-triviale, cette suite est exacte et scindée en vertu de (20.5.12,
(iii)). La commutativité du carré de droite de (20.6.9.1) résulte de la
relation $v=p\circ q$ (20.5.2.7); l'image par
$q_{E/B/A}\otimes1_C$ du noyau $\Upsilon_{C/B/A}$ de $v_{C/B/A}$ est
donc contenue dans le noyau $L$ de $p_{C/E/A}$. D'autre part soit
$h:C\to E$ un $A$-homomorphisme inverse à droite de $p$, et soit
$j:L\to E$ l'injection canonique, de sorte que l'on a
$q(b)=h(v(b))+j(D(b))$ pour $b\in B$, où $D$ est la $A$-dérivation de
$B$ dans $L$ définissant la $B$-extension $E$; on peut écrire
$D=f\circ d_{B/A}$, où $f:\Omega^1_{B/A}\to L$ est un
$B$-homomorphisme. En vertu de (20.5.2.6), on a
\[
  q_{E/B/A}(d_{B/A}(b))=d_{E/A}(q(b))
  =d_{E/A}(h(v(b)))+d_{E/A}(j(f(d_{B/A}(b))))
\]
pour $b\in B$. Soit alors
$z=\sum_i d_{B/A}(b_i)\otimes c_i$, où $b_i\in B$ et $c_i\in C$, un
élément de $\Omega^1_{B/A}\otimes_B C$; on a
\[
  (q_{E/B/A}\otimes1_C)(z)=
  \sum_i d_{E/A}(h(v(b_i)))\otimes c_i+
  \sum_i d_{E/A}(j(f(d_{B/A}(b_i))))\otimes c_i.
\]
Dans la première somme, on a
$d_{E/A}(h(v(b_i)))=h_{E/C/A}(d_{C/A}(v(b_i)))$, donc cette somme est
$(h_{E/C/A}\otimes1_C)(v_{C/B/A}(z))$ en vertu de (20.5.2.6). Si l'on
prend $z\in\Upsilon_{C/B/A}$, cette somme est donc nulle, et il reste,
par définition de $\delta_{C/E/A}$,
\[
  (q_{E/B/A}\otimes1_C)(z)=
  \sum_i\delta_{C/E/A}(c_i f(d_{B/A}(b_i)))
  =\delta_{C/E/A}(f(z)),
\]
ce qui démontre la commutativité du carré de gauche dans (20.6.9.1).
\end{proof}

On notera que lorsque l'on suppose seulement que $\delta_{C/E/A}$ est
\emph{injectif} (et non nécessairement inversible à gauche), cette
interprétation permettrait encore de définir $\chi_E$ comme la
restriction de $q_{E/B/A}\otimes1_C$ à $\Upsilon_{C/B/A}$.

\begin{corollary}[20.6.10]
\label{0.20.6.10-fr}
Si $\mathfrak{R}$ est un idéal de $B$, $C=B/\mathfrak{R}$,
$E=B/\mathfrak{R}^2$, et si $E$ est une $B$-extension $A$-triviale
(18.3.7) de $C$ par $\mathfrak{R}/\mathfrak{R}^2$, alors
l'homomorphisme caractéristique
$\chi_E:\Upsilon_{C/B/A}\to\mathfrak{R}/\mathfrak{R}^2$ est bijectif.
\end{corollary}

\begin{proof}
\oldpage[0\textsubscript{IV-1}]{140}
En effet, dans le diagramme (20.6.9.1), les deux flèches verticales de
droite sont des homomorphismes bijectifs (20.5.12, (ii)).
\end{proof}

\begin{theorem}[20.6.11]
\label{0.20.6.11-fr}
Soient $u:A\to B$, $v:B\to C$ deux homomorphismes d'anneaux, $L$ un
$C$-module. Supposons vérifiée l'une des conditions suivantes :
\begin{enumerate}
  \item[\rm(i)] $L$ est un $C$-module injectif.
  \item[\rm(ii)] $\Upsilon_{C/B/A}$ est facteur direct du $C$-module
$\Omega^1_{B/A}\otimes_B C$ et
$u_{C/B/A}:\Omega^1_{C/A}\to\Omega^1_{C/B}$ est inversible à droite.
\end{enumerate}
Alors l'homomorphisme canonique (20.6.7.1)
\[
  \operatorname{Exalcom}_{B/A}(C,L)\longrightarrow
  \operatorname{Hom}_C(\Upsilon_{C/B/A},L)
\]
est bijectif.

En particulier, si $\Omega^1_{C/B}$ et $\Omega^1_{C/A}$ sont des
$C$-modules projectifs, l'homomorphisme canonique (20.6.7.1) est
bijectif.
\end{theorem}

\begin{proof}
Le fait que chacune des conditions (i), (ii) entraîne que (20.6.7.1) est
bijectif résulte dans les deux cas de la définition de $\alpha_1$. On
notera d'ailleurs que la condition (ii) est \emph{nécessaire et
suffisante} pour que l'homomorphisme (20.6.7.1) soit bijectif pour
\emph{tout} $C$-module $L$ (Bourbaki, \emph{Alg.}, chap. II,
3\textsuperscript{e} éd., \S~2, n\textsuperscript{o}~1, prop.~1). Si on
suppose que $\Omega^1_{C/B}$ et $\Omega^1_{C/A}$ sont des $C$-modules
projectifs, alors, dans la suite exacte (20.6.1.1),
$\operatorname{Ker}(u_{C/B/A})$ est un $C$-module projectif, car la suite exacte
\[
  0\longrightarrow\operatorname{Ker}(u_{C/B/A})\longrightarrow\Omega^1_{C/A}
  \longrightarrow\Omega^1_{C/B}\longrightarrow0
\]
est scindée, $\Omega^1_{C/B}$ étant projectif; comme
$\operatorname{Ker}(u_{C/B/A})=\operatorname{Im}(v_{C/B/A})$, la suite exacte
\[
  0\longrightarrow\Upsilon_{C/B/A}\longrightarrow
  \Omega^1_{B/A}\otimes_B C\longrightarrow
  \operatorname{Im}(v_{C/B/A})\longrightarrow0
\]
est scindée.
\end{proof}

\begin{corollary}[20.6.12]
\label{0.20.6.12-fr}
Supposons que $C$ soit une $A$-algèbre formellement lisse. Il existe alors
un homomorphisme canonique
\begin{equation}
\tag{20.6.12.1}
\label{0.20.6.12.1-fr}
  \operatorname{Exalcom}_B(C,L)\longrightarrow
  \operatorname{Hom}_C(\Upsilon_{C/B/A},L).
\end{equation}
En outre, cet homomorphisme est bijectif si l'une des conditions (i),
(ii) de (20.6.11) est vérifiée.
\end{corollary}

\begin{proof}
En effet, il résulte de l'hypothèse sur $C$ que
$\operatorname{Exalcom}_B(C,L)=\operatorname{Exalcom}_{B/A}(C,L)$, et
l'homomorphisme (20.6.12.1) n'est autre que (20.6.7.1).
\end{proof}

\begin{corollary}[20.6.13]
\label{0.20.6.13-fr}
Si $C$ est une $A$-algèbre formellement lisse et si $\Omega^1_{C/B}$ est
un $C$-module projectif, l'homomorphisme (20.6.12.1) est bijectif.
\end{corollary}

\begin{proof}
En effet, on sait alors que $\Omega^1_{C/A}$ est un $C$-module projectif
(20.4.9), les topologies étant discrètes.
\end{proof}

\begin{env}[20.6.14]
\label{0.20.6.14-fr}
Les notations restant les mêmes, supposons maintenant en outre que $A$,
$B$, $C$ soient des $\Lambda$-algèbres et $u$, $v$ des
$\Lambda$-homomorphismes, ce qui revient à se donner trois homomorphismes
d'anneaux
\[
  \Lambda\xrightarrow{s}A\xrightarrow{u}B\xrightarrow{v}C.
\]

\oldpage[0\textsubscript{IV-1}]{141}
On a donc, en dehors du module d'imperfection $\Upsilon_{C/B/A}$, les
modules d'imperfection $\Upsilon_{B/A/\Lambda}$,
$\Upsilon_{C/A/\Lambda}$ et $\Upsilon_{C/B/\Lambda}$, et on a déjà défini
des homomorphismes canoniques de $C$-modules (20.6.4.2)
\begin{align}
\tag{20.6.14.1}
\label{0.20.6.14.1-fr}
  u':\Upsilon_{C/A/\Lambda}&\longrightarrow\Upsilon_{C/B/\Lambda},\\
\tag{20.6.14.2}
\label{0.20.6.14.2-fr}
  s':\Upsilon_{C/B/\Lambda}&\longrightarrow\Upsilon_{C/B/A}.
\end{align}

Comme dans le diagramme commutatif (20.5.7.3)
\begin{equation}
\tag{20.6.14.3}
\label{0.20.6.14.3-fr}
\xymatrix{
  \Omega^1_{A/\Lambda}\otimes_A B \ar[r] \ar[d] &
  \Omega^1_{B/\Lambda} \ar[r] \ar[d] &
  \Omega^1_{B/A} \ar[r] \ar[d] & 0 \\
  \Omega^1_{A/\Lambda}\otimes_A C \ar[r] &
  \Omega^1_{C/\Lambda} \ar[r] & \Omega^1_{C/A} \ar[r] & 0
}
\end{equation}
la ligne inférieure est formée de $C$-modules, on en déduit par
tensorisation un diagramme commutatif
\begin{equation}
\tag{20.6.14.4}
\label{0.20.6.14.4-fr}
\xymatrix{
  \Omega^1_{A/\Lambda}\otimes_A C \ar[r] \ar@{=}[d] &
  \Omega^1_{B/\Lambda}\otimes_B C \ar[r] \ar[d] &
  \Omega^1_{B/A}\otimes_B C \ar[r] \ar[d] & 0 \\
  \Omega^1_{A/\Lambda}\otimes_A C \ar[r] &
  \Omega^1_{C/\Lambda} \ar[r] & \Omega^1_{C/A} \ar[r] & 0
}
\end{equation}
où la première ligne est encore exacte et la flèche verticale de gauche
est l'identité; si l'on pose
\begin{equation}
\tag{20.6.14.5}
\label{0.20.6.14.5-fr}
  \Upsilon^C_{B/A/\Lambda}=
  \operatorname{Ker}(u_{B/A/\Lambda}\otimes1_C)=
  \operatorname{Ker}(\Omega^1_{A/\Lambda}\otimes_A C\longrightarrow
       \Omega^1_{B/\Lambda}\otimes_B C)
\end{equation}
on voit, compte tenu de la définition de $\Upsilon_{C/A/\Lambda}$, que
l'on a un unique $C$-homomorphisme
\begin{equation}
\tag{20.6.14.6}
\label{0.20.6.14.6-fr}
  v':\Upsilon^C_{B/A/\Lambda}\longrightarrow\Upsilon_{C/A/\Lambda}
\end{equation}
rendant commutatif le diagramme
\begin{equation}
\tag{20.6.14.7}
\label{0.20.6.14.7-fr}
\xymatrix{
  0 \ar[r] & \Upsilon^C_{B/A/\Lambda} \ar[r] \ar[d]^{v'} &
  \Omega^1_{A/\Lambda}\otimes_A C \ar[r] \ar@{=}[d] &
  \Omega^1_{B/\Lambda}\otimes_B C \ar[r] \ar[d] &
  \Omega^1_{B/A}\otimes_B C \ar[r] \ar[d] & 0 \\
  0 \ar[r] & \Upsilon_{C/A/\Lambda} \ar[r] &
  \Omega^1_{A/\Lambda}\otimes_A C \ar[r] &
  \Omega^1_{C/\Lambda} \ar[r] & \Omega^1_{C/A} \ar[r] & 0
}
\end{equation}
dont les lignes sont exactes.

Lorsque $\Lambda=\mathbb{Z}$, on écrira $\Upsilon^C_{B/A}$ au lieu de
$\Upsilon^C_{B/A/\mathbb{Z}}$. Si $\Lambda$ est un corps premier, on a
$\Upsilon^C_{B/A/\Lambda}=\Upsilon^C_{B/A}$.
\end{env}

\begin{env}[20.6.15]
\label{0.20.6.15-fr}
Pour étudier les relations entre les modules précédents, nous introduirons
d'une part le complexe de $B$-modules
$\mathrm{K}_\bullet(B/A/\Lambda)$, d'autre part, les complexes
\oldpage[0\textsubscript{IV-1}]{142}
de $C$-modules $\mathrm{K}_\bullet(C/A/\Lambda)$ et
$\mathrm{K}_\bullet(C/B/\Lambda)$ (20.6.5), et en outre les complexes de
$C$-modules suivants. On posera tout d'abord
\begin{equation}
\tag{20.6.15.1}
\label{0.20.6.15.1-fr}
  \mathrm{K}^C_\bullet(B/A/\Lambda)=
  \mathrm{K}_\bullet(B/A/\Lambda)\otimes_B C.
\end{equation}
D'autre part, nous désignerons par $\mathrm{T}_\bullet(C/B/\Lambda)$ le
complexe de $C$-modules dont les termes sont nuls sauf en degrés 0 et 1,
et où
\begin{equation}
\tag{20.6.15.2}
\label{0.20.6.15.2-fr}
  \mathrm{T}_0(C/B/\Lambda)=\mathrm{T}_1(C/B/\Lambda)=
  \Omega^1_{B/\Lambda}\otimes_B C
\end{equation}
l'opérateur différentiel étant l'identité, de sorte que ce complexe est
homotope à 0; posons enfin
\begin{equation}
\tag{20.6.15.3}
\label{0.20.6.15.3-fr}
  \mathrm{K}'_\bullet(C/A/\Lambda)=
  \mathrm{K}_\bullet(C/A/\Lambda)\oplus
  \mathrm{T}_\bullet(C/B/\Lambda).
\end{equation}
En vertu du caractère trivial de $\mathrm{T}_\bullet$, il est clair que
l'on a
\begin{equation}
\tag{20.6.15.4}
\label{0.20.6.15.4-fr}
  \mathrm{H}^i(\mathrm{K}'_\bullet,L)=
  \mathrm{H}^i(\mathrm{K}_\bullet,L)
  \quad\text{et}\quad
  \mathrm{H}_i(\mathrm{K}'_\bullet,L)=
  \mathrm{H}_i(\mathrm{K}_\bullet,L)
\end{equation}
pour tout $C$-module $L$ et tout $i$.
\end{env}
\begin{env}[20.6.16]
\label{0.20.6.16-fr}
Définissons maintenant une suite exacte de complexes, \emph{scindée} en chaque
degré
\begin{equation}
\tag{20.6.16.1}
\label{0.20.6.16.1-fr}
  0\to\mathrm{K}^C_\bullet(B/A/\Lambda)
  \xrightarrow{j}\mathrm{K}'_\bullet(C/A/\Lambda)
  \xrightarrow{p}\mathrm{K}_\bullet(C/B/\Lambda)\to0
\end{equation}
de la façon suivante : désignons pour un moment par
\[
  f:\Omega^1_{A/\Lambda}\otimes_A C\to
    \Omega^1_{B/\Lambda}\otimes_B C,
  \qquad
  g:\Omega^1_{B/\Lambda}\otimes_B C\to\Omega^1_{C/\Lambda}
\]
les homomorphismes canoniques $u_{B/A/\Lambda}\otimes1_C$ et
$v_{C/B/\Lambda}$ respectivement, dont le composé est
$g\circ f=(v\circ u)_{C/A/\Lambda}$ (cf.\ \sref{0.20.6.14.4-fr}).
On prendra $j_1(x)=(x,f(x))$, $p_1(y,z)=z-f(y)$,
$j_0(x)=(g(x),x)$, $p_0(y,z)=g(z)-y$ de sorte que
$\operatorname{Im}(j_1)=\operatorname{Ker}(p_1)$ est le graphe de $f$,
supplémentaire de $\{0\}\oplus\mathrm{T}_1$, et
$\operatorname{Im}(j_0)=\operatorname{Ker}(p_0)$ est le graphe de $g$,
supplémentaire de $\mathrm{K}_1(C/A/\Lambda)\oplus\{0\}$ ; la vérification
de la commutativité du diagramme
\[
\xymatrix{
  0 \ar[r] & \mathrm{K}^C_1(B/A/\Lambda) \ar[r]^{j_1} \ar[d] &
  \mathrm{K}'_1(C/A/\Lambda) \ar[r]^{p_1} \ar[d] &
  \mathrm{K}_1(C/B/\Lambda) \ar[r] \ar[d] & 0 \\
  0 \ar[r] & \mathrm{K}^C_0(B/A/\Lambda) \ar[r]_{j_0} &
  \mathrm{K}'_0(C/A/\Lambda) \ar[r]_{p_0} &
  \mathrm{K}_0(C/B/\Lambda) \ar[r] & 0
}
\]
où les flèches verticales sont les opérateurs différentiels, est immédiate.
\end{env}

\begin{theorem}[20.6.17]
\label{0.20.6.17-fr}
On a une suite exacte de $C$-modules
\begin{equation}
\tag{20.6.17.1}
\label{0.20.6.17.1-fr}
  0\to\Upsilon^C_{B/A/\Lambda}\xrightarrow{v'}
  \Upsilon_{C/A/\Lambda}\xrightarrow{u'}
  \Upsilon_{C/B/\Lambda}\xrightarrow{\partial}
  \Omega^1_{B/A}\otimes_B C\xrightarrow{v_{C/B/A}}
  \Omega^1_{C/A}\xrightarrow{u_{C/B/A}}\Omega^1_{C/B}\to0
\end{equation}
\oldpage[0\textsubscript{IV-1}]{143}
où l'opérateur bord $\partial$ est le composé
\begin{equation}
\tag{20.6.17.2}
\label{0.20.6.17.2-fr}
  \Upsilon_{C/B/\Lambda}\xrightarrow{s'}\Upsilon_{C/B/A}
  \longrightarrow\Omega^1_{B/A}\otimes_B C
\end{equation}
la seconde flèche étant l'injection canonique.
\end{theorem}

\begin{proof}
En écrivant la suite exacte d'homologie pour la suite exacte de complexes
\sref{0.20.6.16.1-fr}, on obtient \sref{0.20.6.17.1-fr}, l'homologie étant
nulle en degrés distincts de 0 et de 1 ; le fait que les homomorphismes de
cette suite exacte qui proviennent par fonctorialité de $j$ et de $p$ sont
bien ceux de l'énoncé est immédiat. Reste à vérifier que $\partial$ est égal
à \sref{0.20.6.17.2-fr} ; or un élément $z\in\Upsilon_{C/B/\Lambda}$ est
image par $p_1$ de $(0,z)$, d'où on déduit aussitôt que $\partial(z)$ est
l'image de $z$ par l'homomorphisme canonique
\[
  s_{B/A/\Lambda}\otimes1_C:
  \Omega^1_{B/\Lambda}\otimes_B C\to\Omega^1_{B/A}\otimes_B C.
\]
Notre assertion résulte de la commutativité du diagramme
\begin{equation}
\tag{20.6.17.3}
\label{0.20.6.17.3-fr}
\xymatrix{
  \Upsilon_{C/B/\Lambda} \ar[r]^{s'} \ar[dd] & \Upsilon_{C/B/A} \ar[dd] \\
  & \\
  \Omega^1_{B/\Lambda}\otimes_B C \ar[r] &
  \Omega^1_{B/A}\otimes_B C
}
\end{equation}
(cf.\ \sref{0.20.6.4.3-fr}).
\end{proof}

\begin{corollary}[20.6.18]
\label{0.20.6.18-fr}
\begin{enumerate}
  \item[\rm(i)] La suite de $C$-modules
\begin{equation}
\tag{20.6.18.1}
\label{0.20.6.18.1-fr}
  0\to\Upsilon^C_{B/A/\Lambda}\xrightarrow{v'}
  \Upsilon_{C/A/\Lambda}\xrightarrow{u'}
  \Upsilon_{C/B/\Lambda}\xrightarrow{s'}\Upsilon_{C/B/A}\to0
\end{equation}
est exacte.
  \item[\rm(ii)] Si $B$ est une $A$-algèbre formellement lisse relativement
  à $\Lambda$, on a $\Upsilon^C_{B/A/\Lambda}=0$.
\end{enumerate}
\end{corollary}

\begin{proof}
(i) Dans \sref{0.20.6.17.1-fr}, l'image de $\Upsilon_{C/B/\Lambda}$ par
$\partial$ est le noyau de $v_{C/B/A}$, donc $\Upsilon_{C/B/A}$ par
définition.

(ii) L'hypothèse entraîne que la suite
\[
  0\to\Omega^1_{A/\Lambda}\otimes_A B\to
  \Omega^1_{B/\Lambda}\to\Omega^1_{B/A}\to0
\]
est exacte et scindée \sref{0.20.5.7-fr} ; par tensorisation avec $C$, la
suite
\[
  0\to\Omega^1_{A/\Lambda}\otimes_A C\to
  \Omega^1_{B/\Lambda}\otimes_B C\to
  \Omega^1_{B/A}\otimes_B C\to0
\]
reste donc exacte, d'où notre assertion.
\end{proof}

\begin{corollary}[20.6.19]
\label{0.20.6.19-fr}
Soient $K$ un corps, $E$, $F$ deux extensions de $K$ telles que
$K\subset E\subset F$.
\begin{enumerate}
  \item[\rm(i)] Si $F$ est une extension séparable de $E$, on a
  $\Upsilon_{F/E/K}=0$ (autrement dit, l'homomorphisme canonique
  $\Omega^1_{E/K}\otimes_E F\to\Omega^1_{F/K}$ est injectif).
  \item[\rm(ii)] Inversement, si $F$ est une extension séparable de $K$ et
  si l'on a $\Upsilon_{F/E/K}=0$, alors $F$ est une extension séparable de
  $E$.
\end{enumerate}
\end{corollary}

\begin{proof}
Si $P$ est le corps premier de $K$, on a en effet la suite exacte
\sref{0.20.6.18-fr}
\[
  \Upsilon_{F/K/P}\to\Upsilon_{F/E/P}\to\Upsilon_{F/E/K}\to0.
\]
Si $\Upsilon_{F/E/P}=\Upsilon_{F/E}=0$, on a donc
$\Upsilon_{F/E/K}=0$, d'où (i) en vertu de \sref{0.20.6.3-fr} ;
inversement, si $\Upsilon_{F/E/K}=0$ et
$\Upsilon_{F/K/P}=\Upsilon_{F/K}=0$, on a
$\Upsilon_{F/E/P}=0$, d'où (ii) en vertu de \sref{0.20.6.3-fr}.
\end{proof}

\oldpage[0\textsubscript{IV-1}]{144}
\begin{corollary}[20.6.20]
\label{0.20.6.20-fr}
\begin{enumerate}
  \item[\rm(i)] Si $K$ est une extension algébrique séparable de $k$, on a
  $\Omega^1_{K/k}=0$.
  \item[\rm(ii)] Soient $k$ un corps de caractéristique 0, $K$ une extension
  de $k$. Pour que $\Omega^1_{K/k}=0$, il faut et il suffit que $K$ soit une
  extension algébrique de $k$. En particulier, pour qu'un corps $K$ de
  caractéristique 0 soit tel que $\Omega^1_K=0$, il faut et il suffit que
  $K$ soit une extension algébrique de $\mathbb{Q}$ (cf.\ (21.4.4) et
  (21.7.5)).
\end{enumerate}
\end{corollary}

\begin{proof}
(i) Pour tout $x\in K$, soit $f$ le polynôme minimal de $x$ sur $k$ ; comme
$f'(x)\ne0$ et $d_{K/k}(f(x))=f'(x)d_{K/k}(x)=0$, on a
$d_{K/k}(x)=0$ et notre assertion résulte de \sref{0.20.4.7-fr}.

(ii) Il existe une extension pure $L$ de $k$ telle que
$k\subset L\subset K$ et que $K$ soit une extension algébrique de $L$.
Comme $K$ est séparable sur $L$, il résulte de \sref{0.20.6.19-fr}, (i) que
la suite \sref{0.20.5.7.2-fr}
\[
  0\to\Omega^1_{L/k}\otimes_L K\to\Omega^1_{K/k}
  \to\Omega^1_{K/L}\to0
\]
est exacte et de (i) que $\Omega^1_{K/L}=0$. La relation
$\Omega^1_{K/k}=0$ est donc équivalente à $\Omega^1_{L/k}=0$, et comme
$L$ est une extension pure de $k$, il résulte de \sref{0.20.5.10-fr} que
la relation $\Omega^1_{L/k}=0$ équivaut à $L=k$.
\end{proof}

\begin{remarks}[20.6.21]
\label{0.20.6.21-fr}
\begin{enumerate}
  \item[\rm(i)] Comme $\Upsilon_{B/A/\Lambda}$ est le noyau de
  \[
    u_{B/A/\Lambda}:\Omega^1_{A/\Lambda}\otimes_A B
    \to\Omega^1_{B/\Lambda}
  \]
  et $\Upsilon^C_{B/A/\Lambda}$ le noyau de
  $u_{B/A/\Lambda}\otimes1_C$, on a un homomorphisme canonique
\begin{equation}
\tag{20.6.21.1}
\label{0.20.6.21.1-fr}
  \Upsilon_{B/A/\Lambda}\otimes_B C\to\Upsilon^C_{B/A/\Lambda}.
\end{equation}
Cet homomorphisme est bijectif lorsque la suite
\begin{equation}
\tag{20.6.21.2}
\label{0.20.6.21.2-fr}
  0\to\Upsilon_{B/A/\Lambda}\otimes_B C\to
  \Omega^1_{A/\Lambda}\otimes_A C\to
  \Omega^1_{B/\Lambda}\otimes_B C\to
  \Omega^1_{B/A}\otimes_B C\to0
\end{equation}
est exacte, ce qui se produit dans les cas suivants :
\begin{enumerate}
  \item[$1^\circ$] $C$ est un $B$-module \emph{plat}.
  \item[$2^\circ$] Les $B$-modules $\Omega^1_{B/\Lambda}$ et
  $\Omega^1_{B/A}$ sont \emph{plats} : en effet, il en est alors de même de
  $\operatorname{Ker}(\Omega^1_{B/\Lambda}\to\Omega^1_{B/A})$
  ($\mathbf{0_I}$, 6.1.2), et la suite \sref{0.20.6.21.2-fr} est alors
  exacte en vertu de ($\mathbf{0_I}$, 6.1.2).
\end{enumerate}
  \item[\rm(ii)] Considérons un diagramme commutatif d'homomorphismes
  d'anneaux
\[
\xymatrix{
  \Lambda' \ar[r] & A' \ar[r] & B' \ar[r] & C' \\
  \Lambda \ar[u] \ar[r] & A \ar[u] \ar[r] & B \ar[u] \ar[r] & C \ar[u]
}
\]
Alors les définitions de \sref{0.20.6.16-fr} montrent que l'on a un
diagramme commutatif de complexes (où les flèches verticales proviennent de
\sref{0.20.6.6.4-fr})
\[
\xymatrix{
  0 \ar[r] & \mathrm{K}^C_\bullet(B/A/\Lambda) \ar[r] \ar[d] &
  \mathrm{K}'_\bullet(C/A/\Lambda) \ar[r] \ar[d] &
  \mathrm{K}_\bullet(C/B/\Lambda) \ar[r] \ar[d] & 0 \\
  0 \ar[r] & \mathrm{K}^{C'}_\bullet(B'/A'/\Lambda') \ar[r] &
  \mathrm{K}'_\bullet(C'/A'/\Lambda') \ar[r] &
  \mathrm{K}_\bullet(C'/B'/\Lambda') \ar[r] & 0
}
\]
\oldpage[0\textsubscript{IV-1}]{145}
d'où, par passage à l'homologie, un \emph{diagramme commutatif}
\begin{equation}
\tag{20.6.21.3}
\label{0.20.6.21.3-fr}
\xymatrix@C=0.7em{
  0 \ar[r] & \Upsilon^C_{B/A/\Lambda} \ar[r] \ar[d] &
  \Upsilon_{C/A/\Lambda} \ar[r] \ar[d] &
  \Upsilon_{C/B/\Lambda} \ar[r] \ar[d] &
  \Omega^1_{B/A}\otimes_B C \ar[r] \ar[d] &
  \Omega^1_{C/A} \ar[r] \ar[d] &
  \Omega^1_{C/B} \ar[r] \ar[d] & 0 \\
  0 \ar[r] & \Upsilon^{C'}_{B'/A'/\Lambda'} \ar[r] &
  \Upsilon_{C'/A'/\Lambda'} \ar[r] &
  \Upsilon_{C'/B'/\Lambda'} \ar[r] &
  \Omega^1_{B'/A'}\otimes_{B'} C' \ar[r] &
  \Omega^1_{C'/A'} \ar[r] & \Omega^1_{C'/B'} \ar[r] & 0
}
\end{equation}
On a un diagramme commutatif analogue pour \sref{0.20.6.18.1-fr}.
\end{enumerate}
\end{remarks}

\begin{proposition}[20.6.22]
\label{0.20.6.22-fr}
Soient $s:\Lambda\to A$, $u:A\to B$ deux homomorphismes d'anneaux,
$\mathfrak{R}$ un idéal de $B$, $C$ l'anneau quotient $B/\mathfrak{R}$.
Supposons que $E=B/\mathfrak{R}^2$ soit une $B$-extension
$\Lambda$-triviale de $C$ par $\mathfrak{R}/\mathfrak{R}^2$. On a alors une
suite exacte
\begin{equation}
\tag{20.6.22.1}
\label{0.20.6.22.1-fr}
  0\to\Upsilon^C_{B/A/\Lambda}\xrightarrow{v'}
  \Upsilon_{C/A/\Lambda}\xrightarrow{\chi_E}
  \mathfrak{R}/\mathfrak{R}^2\xrightarrow{\delta_{C/B/A}}
  \Omega^1_{B/A}\otimes_B C\xrightarrow{v_{C/B/A}}
  \Omega^1_{C/A}\to0.
\end{equation}
\end{proposition}

\begin{proof}
En effet, comme $v:B\to C$ est surjectif, on a $\Omega^1_{C/B}=0$
\sref{0.20.4.12-fr}. En outre, il résulte de \sref{0.20.6.10-fr} que
$\Upsilon_{C/B/\Lambda}$ s'identifie canoniquement à
$\mathfrak{R}/\mathfrak{R}^2$. Il suffit alors d'appliquer les suites exactes
\sref{0.20.6.17.1-fr} et \sref{0.20.6.18.1-fr}, en notant que l'on a un
diagramme commutatif
\[
\xymatrix{
  \Upsilon_{C/A/\Lambda} \ar[r] \ar[d]_{\chi_E} &
  \Omega^1_{A/\Lambda}\otimes_A C \ar[d] \\
  \Upsilon_{C/B/\Lambda}=\Upsilon_{C/B/A}=\mathfrak{R}/\mathfrak{R}^2
    \ar[r]^-{\delta_{C/B/A}} \ar[d]_{\ell} &
  \Omega^1_{B/A}\otimes_B C=\Omega^1_{E/A}\otimes_E C \ar[d] \\
  \mathfrak{R}/\mathfrak{R}^2 \ar[r]^{\delta_{C/B/A}} &
  \Omega^1_{B/A}\otimes_B C
}
\]
et utilisant la commutativité du diagramme \sref{0.20.6.17.3-fr}.
\end{proof}

On a ainsi précisé le \emph{noyau} de $\delta_{C/B/A}$ dans le cas où il
existe un anneau $\Lambda$ tel que $A$ soit une $\Lambda$-algèbre et que
$B/\mathfrak{R}^2$ soit $\Lambda$-triviale ; on observera que c'est le cas
en particulier lorsque $C$ est une \emph{$\Lambda$-algèbre formellement
lisse} (pour la topologie discrète).

Supposons en outre qu'on ait un diagramme commutatif d'homomorphismes
d'anneaux
\[
\xymatrix{
  \Lambda' \ar[r] & A' \ar[r] & B' \\
  \Lambda \ar[u] \ar[r] & A \ar[u] \ar[r] & B \ar[u]^f
}
\]
\oldpage[0\textsubscript{IV-1}]{146}
que $\mathfrak{R}'$ soit un idéal de $B'$ tel que
$f(\mathfrak{R})\subset\mathfrak{R}'$, et que
$E'=B'/\mathfrak{R}'^2$ soit une $B'$-extension $\Lambda'$-triviale de
$C'=B'/\mathfrak{R}'$ par $\mathfrak{R}'/\mathfrak{R}'^2$. On a alors un
diagramme commutatif
\begin{equation}
\tag{20.6.22.2}
\label{0.20.6.22.2-fr}
\xymatrix@C=0.7em{
  0 \ar[r] & \Upsilon^{C'}_{B'/A'/\Lambda'} \ar[r] &
  \Upsilon_{C'/A'/\Lambda'} \ar[r] &
  \mathfrak{R}'/\mathfrak{R}'^2 \ar[r] &
  \Omega^1_{B'/A'}\otimes_{B'} C' \ar[r] &
  \Omega^1_{C'/A'} \ar[r] & 0 \\
  0 \ar[r] & \Upsilon^C_{B/A/\Lambda} \ar[r] \ar[u] &
  \Upsilon_{C/A/\Lambda} \ar[r] \ar[u] &
  \mathfrak{R}/\mathfrak{R}^2 \ar[r] \ar[u] &
  \Omega^1_{B/A}\otimes_B C \ar[r] \ar[u] &
  \Omega^1_{C/A} \ar[r] \ar[u] & 0
}
\end{equation}
comme il résulte de \sref{0.20.6.21.3-fr} et \sref{0.20.5.11.3-fr}.

\begin{corollary}[20.6.23]
\label{0.20.6.23-fr}
Sous les hypothèses de \sref{0.20.6.22-fr}, supposons en outre que $B$ soit
une $A$-algèbre formellement lisse. Alors on a une suite exacte
\begin{equation}
\tag{20.6.23.1}
\label{0.20.6.23.1-fr}
  0\to\Upsilon_{C/A/\Lambda}\xrightarrow{\chi_E}
  \mathfrak{R}/\mathfrak{R}^2\xrightarrow{\delta_{C/B/A}}
  \Omega^1_{B/A}\otimes_B C\to\Omega^1_{C/A}\to0.
\end{equation}
\end{corollary}

\begin{proof}
Cela résulte en effet de \sref{0.20.6.18-fr}, (ii).
\end{proof}

\begin{env}[20.6.24]
\label{0.20.6.24-fr}
Lorsque les hypothèses de \sref{0.20.6.22-fr} sont satisfaites, on dit encore
que l'homomorphisme caractéristique $\chi_E$ est
l'\emph{homomorphisme caractéristique de la $A$-algèbre $B$, relativement à
$\Lambda$ et à l'idéal $\mathfrak{R}$} (en supprimant ces dernières
précisions lorsqu'il n'y a pas de confusion à craindre) ; on le notera
parfois $\chi_B$ ou $\chi_{B/A}$.
\end{env}

\begin{proposition}[20.6.25]
\label{0.20.6.25-fr}
Soient $s:\Lambda\to A$, $u:A\to B$, $v:B\to C$ trois homomorphismes
d'anneaux, $L$ un $C$-module. On a alors une suite exacte
\begin{equation}
\tag{20.6.25.1}
\label{0.20.6.25.1-fr}
\begin{split}
  0\to\operatorname{D\acute{e}r}_B(C,L)\xrightarrow{u^1}
  \operatorname{D\acute{e}r}_A(C,L)\xrightarrow{v^1}
  \operatorname{D\acute{e}r}_A(B,L)\xrightarrow{\partial} \\
  \to\operatorname{Exalcom}_{B/\Lambda}(C,L)\xrightarrow{u^1}
  \operatorname{Exalcom}_{A/\Lambda}(C,L)\xrightarrow{v^1}
  \operatorname{Exalcom}_{A/\Lambda}(B,L)\to0
\end{split}
\end{equation}
où $u^1$, $v^1$ sont les homomorphismes définis dans (18.3.6.4) et
(18.3.6.2), et $\partial$ est défini comme dans \sref{0.20.2.2-fr}.
\end{proposition}

\begin{proof}
En effet, comme la suite exacte \sref{0.20.6.16.1-fr} est scindée, on en
déduit une suite exacte
\[
  0\to\operatorname{Hom}_C(\mathrm{K}_\bullet(C/B/\Lambda),L)\to
  \operatorname{Hom}_C(\mathrm{K}'_\bullet(C/A/\Lambda),L)\to
  \operatorname{Hom}_C(\mathrm{K}^C_\bullet(B/A/\Lambda),L)\to0.
\]
Si on applique à ce complexe la suite exacte de cohomologie, en tenant compte
de \sref{0.20.6.15.4-fr}, et de \sref{0.20.6.6-fr}, on obtient
\sref{0.20.6.25.1-fr} car on a
\[
  \operatorname{Hom}_C(\mathrm{K}^C_\bullet(B/A/\Lambda),L)
  =\operatorname{Hom}_B(\mathrm{K}_\bullet(B/A/\Lambda),L)
\]
par définition ; l'identification de $u^1$ et $v^1$ avec les homomorphismes
de (18.3.4.2) résulte de \sref{0.20.6.6.4-fr}.
\end{proof}

\begin{remark}[20.6.26]
\label{0.20.6.26-fr}
Dans ce numéro, les complexes $\mathrm{K}_\bullet(C/B/A)$ sont apparus comme
un artifice technique destiné à simplifier l'exposition de certains
comportements fonctoriels. En réalité, ces complexes, considérés comme objets
de la catégorie des
\oldpage[0\textsubscript{IV-1}]{147}
complexes de $C$-modules « à homotopie près » (c'est-à-dire où les morphismes
sont les classes d'homomorphismes homotopes) sont des invariants remarquables,
plus fins que le couple formé de $\Omega^1_{C/B}$ et de
$\Upsilon_{C/B/A}$. Lorsque $k$ est un corps premier, et $C$ une $k$-algèbre
formellement lisse (par exemple un anneau régulier de type fini sur une
extension de $k$ (cf.\ ($\mathbf{IV}$, 6.8.6))), on peut montrer que le
complexe $\mathrm{K}_\bullet(C/A/k)$ peut se décrire uniquement en faisant
intervenir $C$ et $A$ (à l'exclusion de $k$) : on exprime $C$ comme quotient
d'une algèbre de polynômes $B$ sur $A$ par un idéal $\mathfrak{Q}$, et on
considère le complexe $\mathrm{F}_\bullet(C/A)$ à 2 termes non nuls
\[
  \cdots\to0\to\mathfrak{Q}/\mathfrak{Q}^2\to
  \Omega^1_{B/A}\otimes_B C\to0\to\cdots
\]
(dont l'homologie coïncide bien avec celle de
$\mathrm{K}_\bullet(C/A/k)$ en vertu de \sref{0.20.6.23.1-fr}). Ces
complexes $\mathrm{F}_\bullet(C/A)$, qui du point de vue de l'algèbre
homologique jouent le rôle d'un fibré conormal pour $\operatorname{Spec}(C)$
au-dessus de $\operatorname{Spec}(A)$, tiendront une place importante dans
les chapitres de cet ouvrage consacrés à la dualité des faisceaux cohérents et
au théorème de Riemann--Roch.
\end{remark}

\subsection{Généralisations aux anneaux topologiques}
\label{subsection:0.20.7-fr}

\begin{env}[20.7.1]
\label{0.20.7.1-fr}
Il résulte aussitôt des définitions que si, dans \sref{0.20.5.2-fr} et
\sref{0.20.5.3-fr}, les anneaux $A$, $B$, $C$ sont supposés être des anneaux
topologiques et les homomorphismes d'anneaux $u$, $v$ continus, alors les
homomorphismes $u_{C/B/A}$ et $v_{C/B/A}$ sont \emph{continus} (il faut
naturellement prendre sur $\Omega^1_{B/A}\otimes_B C$ la topologie produit
tensoriel).

En outre, $u_{C/B/A}:\Omega^1_{C/A}\to\Omega^1_{C/B}$ est un
\emph{morphisme strict surjectif} de $C$-modules topologiques ; en effet, il en
est ainsi de l'homomorphisme canonique $C\otimes_A C\to C\otimes_B C$, compte
tenu de la définition de la topologie produit tensoriel ; par passage aux
quotients on en déduit que l'homomorphisme correspondant
$\mathrm{P}^1_{C/A}\to\mathrm{P}^1_{C/B}$ est un morphisme strict surjectif,
et $u_{C/B/A}$ est la restriction de ce dernier à $\Omega^1_{C/A}$.

Dans \sref{0.20.5.5-fr}, si $A'$ et $B$ sont des $A$-algèbres topologiques,
et si $B'$ et $\Omega^1_{B/A}\otimes_B B'$ sont munis des topologies produits
tensoriels, l'homomorphisme canonique \sref{0.20.5.5.1-fr} est un isomorphisme
topologique, en tenant compte de ce que $\mathrm{P}^1_{B/A}$ est somme directe
topologique de $B$ et de $\Omega^1_{B/A}$ \sref{0.20.4.8-fr}.
\end{env}

\begin{proposition}[20.7.2]
\label{0.20.7.2-fr}
Soient $u:A\to B$, $v:B\to C$ deux homomorphismes continus d'anneaux
topologiques. Pour que l'homomorphisme continu
$v_{C/B/A}:\Omega^1_{B/A}\otimes_B C\to\Omega^1_{C/A}$ soit formellement
inversible à gauche (cf.\ \sref{0.19.1.5-fr}), il faut et il suffit que $C$
soit une $B$-algèbre formellement lisse relativement à $A$
\sref{0.19.9.1-fr} (et \emph{a fortiori} il suffit que $C$ soit une
$B$-algèbre formellement lisse).
\end{proposition}

\begin{proof}
Dire que $v_{C/B/A}$ est formellement inversible à gauche signifie en effet,
vu que les topologies de $\Omega^1_{C/A}$ et $\Omega^1_{B/A}\otimes_B C$
sont moins fines que celles déduites de la topologie de $C$
\sref{0.20.4.5-fr}, que pour tout $C$-module \emph{discret} $L$, annulé par un
idéal ouvert de $C$, l'homomorphisme canonique
\[
  \operatorname{Hom.\ cont}_C(\Omega^1_{C/A},L)\to
  \operatorname{Hom.\ cont}_C(\Omega^1_{B/A}\otimes_B C,L)
\]
\oldpage[0\textsubscript{IV-1}]{148}
est \emph{surjectif} \sref{0.19.1.5-fr} ; comme
$\operatorname{Hom.\ cont}_C(\Omega^1_{B/A}\otimes_B C,L)=
\operatorname{Hom.\ cont}_B(\Omega^1_{B/A},L)$ par définition de la topologie
produit tensoriel, il revient au même, en vertu de \sref{0.20.4.8-fr}, de dire
que l'homomorphisme canonique
\[
  \operatorname{D\acute{e}r.\ cont}_A(C,L)\to
  \operatorname{D\acute{e}r.\ cont}_A(B,L)
\]
est surjectif. Mais la suite exacte \sref{0.20.3.8.2-fr} montre que cette
condition équivaut à $\operatorname{Exalcotop}_{B/A}(C,L)=0$, c'est-à-dire
précisément au fait que $C$ est formellement lisse relativement à $A$
\sref{0.19.9.8-fr}.
\end{proof}

\begin{corollary}[20.7.3]
\label{0.20.7.3-fr}
Supposons que dans $B$ et dans $C$, le carré d'un idéal ouvert soit ouvert.
Pour que $C$ soit une $B$-algèbre formellement lisse relativement à $A$, il
faut et il suffit que si l'on désigne par $(\mathfrak{R}_\lambda)$ un système
fondamental de voisinages de $0$ formé d'idéaux de $C$, alors, pour tout
$\lambda$, l'homomorphisme
\begin{equation}
\tag{20.7.3.1}
\label{0.20.7.3.1-fr}
  v_{C/B/A}\otimes 1_{C/\mathfrak{R}_\lambda}:
  \Omega^1_{B/A}\otimes_B(C/\mathfrak{R}_\lambda)\to
  \Omega^1_{C/A}\otimes_C(C/\mathfrak{R}_\lambda)
\end{equation}
soit inversible à gauche.
\end{corollary}

\begin{proof}
On sait en effet alors que la topologie de $\Omega^1_{B/A}$ (resp.\
$\Omega^1_{C/A}$) est déduite de celle de $B$ (resp.\ de $C$)
\sref{0.20.4.5-fr} ; on en conclut aussitôt que la topologie de
$\Omega^1_{B/A}\otimes_B C$ est aussi déduite de celle de $C$ ; le corollaire
résulte alors de \sref{0.20.7.2-fr} et \sref{0.19.1.7-fr}.
\end{proof}

\begin{proposition}[20.7.4]
\label{0.20.7.4-fr}
Soient $A$ un anneau topologique, $B$ une $A$-algèbre topologique. Pour que
$B$ soit formellement non ramifiée \sref{0.19.10.2-fr}, il faut et il suffit
que le séparé complété $\widehat{\Omega}^1_{B/A}$ soit nul.
\end{proposition}

\begin{proof}
En effet, il résulte aussitôt de \sref{0.19.10.4-fr} et
\sref{0.20.1.1-fr} que, pour que $B$ soit formellement non ramifiée, il faut et
il suffit que pour tout idéal ouvert $\mathfrak{R}$ de $B$, tout idéal ouvert
$\mathfrak{H}$ de $A$ tel que $\mathfrak{H}B\subset\mathfrak{R}$ et tout
$(B/\mathfrak{R})$-module $L$, on ait
$\operatorname{D\acute{e}r}_{A/\mathfrak{H}}(B/\mathfrak{R},L)=0$,
c'est-à-dire $\operatorname{D\acute{e}r.\ cont}_A(B,L)=0$ pour tout
$B$-module discret $L$ annulé par un idéal ouvert de $B$ ; en vertu de
\sref{0.20.4.8-fr}, cela équivaut à
$\operatorname{Hom.\ cont}_B(\Omega^1_{B/A},L)=0$ pour un tel $B$-module,
d'où aussitôt la proposition.

Lorsque $B$ est \emph{discret}, la condition de l'énoncé de
\sref{0.20.7.4-fr} équivaut donc à $\Omega^1_{B/A}=0$.
\end{proof}

\begin{corollary}[20.7.5]
\label{0.20.7.5-fr}
Soient $A$ un anneau, $\mathfrak{m}$ un idéal de $A$, $B$ une $A$-algèbre ;
munissons $A$ de la topologie $\mathfrak{m}$-préadique, $B$ de la topologie
déduite de celle de $A$, et posons $A_0=A/\mathfrak{m}$,
$B_0=B/\mathfrak{m}B=B\otimes_A A_0$. Alors, pour que $B$ soit une
$A$-algèbre formellement non ramifiée, il faut et il suffit que
$\Omega^1_{B_0/A_0}=0$ (ou encore que $B_0$ soit une $A_0$-algèbre
formellement non ramifiée).
\end{corollary}

\begin{proof}
En effet, on a
$\Omega^1_{B/A}\otimes_B B_0=\Omega^1_{B/A}/\mathfrak{m}.\Omega^1_{B/A}
=\Omega^1_{B_0/A_0}$ en vertu de \sref{0.20.5.5-fr} ; écrire que tout
sous-module ouvert de $\Omega^1_{B/A}$ est égal à $\Omega^1_{B/A}$ équivaut
par \sref{0.20.4.5-fr} à écrire que
$\mathfrak{m}.\Omega^1_{B/A}=\Omega^1_{B/A}$, d'où la conclusion.
\end{proof}

\begin{proposition}[20.7.6]
\label{0.20.7.6-fr}
Soient $u:A\to B$, $v:B\to C$ deux homomorphismes continus d'anneaux
topologiques, et supposons que $v$ fasse de $C$ une $B$-algèbre formellement
étale ; alors $v_{C/B/A}:\Omega^1_{B/A}\otimes_B C\to\Omega^1_{C/A}$ est un
bimorphisme formel \sref{0.19.1.2-fr}.
\end{proposition}

\begin{proof}
En effet, pour tout $C$-module discret $L$ annulé par un idéal ouvert de $C$,
on a alors \sref{0.20.7.4-fr}
$\operatorname{D\acute{e}r.\ cont}_B(C,L)=0$ et par suite
\sref{0.20.3.6.1-fr} l'homomorphisme canonique
\oldpage[0\textsubscript{IV-1}]{149}
$\operatorname{D\acute{e}r.\ cont}_A(C,L)\to
\operatorname{D\acute{e}r.\ cont}_A(B,L)$ est injectif ; il est par ailleurs
surjectif en vertu de \sref{0.20.7.2-fr}, donc il est bijectif ; il en résulte
que l'image de $v_{C/B/A}$ est nécessairement dense dans $\Omega^1_{C/A}$,
sans quoi le quotient de $\Omega^1_{C/A}$ par l'adhérence de
$\operatorname{Im}(v_{C/B/A})$ serait séparé et $\neq0$ et aurait donc un
quotient discret $L\neq0$, contrairement à ce qu'on vient de voir (compte tenu
de \sref{0.20.4.8-fr}). Par suite $v_{C/B/A}$, qui est un monomorphisme formel
en vertu de \sref{0.20.7.2-fr}, est aussi un épimorphisme formel
\sref{0.19.1.2-fr}, donc un bimorphisme formel.
\end{proof}

\begin{corollary}[20.7.7]
\label{0.20.7.7-fr}
Supposons que dans $B$ et dans $C$ le carré de tout idéal ouvert soit ouvert.
Si $C$ est une $B$-algèbre formellement étale, alors, pour tout idéal ouvert
$\mathfrak{R}_\lambda$ de $C$, l'homomorphisme \sref{0.20.7.3.1-fr} est
bijectif.
\end{corollary}

\begin{proof}
En effet, en vertu de \sref{0.19.1.1-fr}, cet homomorphisme est surjectif, et
il est injectif par \sref{0.20.7.3-fr}.
\end{proof}

\begin{proposition}[20.7.8]
\label{0.20.7.8-fr}
Soient $u:A\to B$ un homomorphisme continu d'anneaux topologiques,
$\mathfrak{R}$ un idéal de $B$, $C$ l'anneau topologique quotient
$B/\mathfrak{R}$, $v:B\to C$ l'homomorphisme canonique. Alors :
\begin{enumerate}
  \item[\rm(i)] Dans la suite exacte \sref{0.20.5.12.1-fr}
\[
  \mathfrak{R}/\mathfrak{R}^2\xrightarrow{\delta_{C/B/A}}
  \Omega^1_{B/A}\otimes_B C\xrightarrow{v_{C/B/A}}\Omega^1_{C/A}\to0
\]
l'homomorphisme $\delta_{C/B/A}$ est continu et l'homomorphisme $v_{C/B/A}$
est un morphisme strict de $C$-modules topologiques.
  \item[\rm(ii)] Pour que $\delta_{C/B/A}$ soit formellement inversible à
gauche \sref{0.19.1.5-fr}, il faut et il suffit que pour tout $C$-module
discret $L$ annulé par un idéal ouvert de $C$, l'homomorphisme canonique
\begin{equation}
\tag{20.7.8.1}
\label{0.20.7.8.1-fr}
  \operatorname{Exalcotop}_A(C,L)\to\operatorname{Exalcotop}_A(B,L)
\end{equation}
soit injectif.
\end{enumerate}
\end{proposition}

\begin{proof}
(i) La première assertion est évidente. D'autre part, l'homomorphisme
canonique $v\otimes v:B\otimes_A B\to C\otimes_A C$ est un morphisme strict
par définition de la topologie produit tensoriel, et on en déduit aussitôt
(cf.\ \sref{0.20.5.2-fr}) qu'il en est de même de $v_{C/B/A}$.

(ii) Dire que $\delta_{C/B/A}$ est formellement inversible à gauche signifie
que pour tout $C$-module discret $L$ annulé par un idéal ouvert de $C$,
l'homomorphisme canonique
\[
  \operatorname{Hom.\ cont}_C(\Omega^1_{B/A}\otimes_B C,L)\to
  \operatorname{Hom.\ cont}_C(\mathfrak{R}/\mathfrak{R}^2,L)
\]
est surjectif. Or, compte tenu de \sref{0.18.4.3-fr} et de
\sref{0.20.4.8-fr}, cela revient à dire que l'homomorphisme canonique
\[
  \operatorname{D\acute{e}r.\ cont}_A(B,L)\to
  \operatorname{Exalcotop}_B(C,L)
\]
est surjectif, et la conclusion résulte donc de la suite exacte
\sref{0.20.3.6.1-fr}.

Le fait que \sref{0.20.7.8.1-fr} soit injectif s'exprime encore de la manière
suivante, compte tenu de la définition des deux membres
\sref{0.18.4.1-fr} : pour tout idéal ouvert $\mathfrak{M}$ de $A$, tout idéal
ouvert $\mathfrak{N}$ de $B$ tel que $\mathfrak{M}B\subset\mathfrak{N}$, et
toute $(A/\mathfrak{M})$-extension $E$ de
$B/(\mathfrak{R}+\mathfrak{N})$ par un
$(B/(\mathfrak{R}+\mathfrak{N}))$-module $L$, telle que l'image réciproque de
$E$ par l'homomorphisme canonique
$B/\mathfrak{N}\to B/(\mathfrak{R}+\mathfrak{N})$ soit
$(A/\mathfrak{M})$-triviale, il existe un idéal ouvert
$\mathfrak{M}'\subset\mathfrak{M}$ de $A$,
\oldpage[0\textsubscript{IV-1}]{150}
un idéal ouvert $\mathfrak{N}'\subset\mathfrak{N}$ de $B$ tels que
$\mathfrak{M}'B\subset\mathfrak{N}'$ et que l'image réciproque de $E$ par
l'homomorphisme canonique
$B/(\mathfrak{R}+\mathfrak{N}')\to B/(\mathfrak{R}+\mathfrak{N})$ soit
$(A/\mathfrak{M}')$-triviale.
\end{proof}

\begin{corollary}[20.7.9]
\label{0.20.7.9-fr}
Si la $A$-algèbre topologique $C=B/\mathfrak{R}$ est formellement lisse,
l'homomorphisme canonique $\delta_{C/B/A}$ est formellement inversible à
gauche.
\end{corollary}

\begin{env}[20.7.10]
\label{0.20.7.10-fr}
Dans \sref{0.20.6.1-fr}, lorsque $A$, $B$, $C$ sont des anneaux topologiques,
$u$ et $v$ des homomorphismes continus, on munit $\Upsilon_{C/B/A}$ de la
topologie induite par celle de $\Omega^1_{B/A}\otimes_B C$ ; les
homomorphismes \sref{0.20.6.4.2-fr} et \sref{0.20.6.4.4-fr} sont alors
continus pourvu qu'il en soit de même de ceux du diagramme
\sref{0.20.6.4.1-fr}. En outre, si, dans \sref{0.20.6.7-fr}, on suppose que
$L$ est un $C$-module discret annulé par un idéal ouvert de $C$, on déduit,
par passage à la limite inductive, un $C$-homomorphisme canonique
\begin{equation}
\tag{20.7.10.1}
\label{0.20.7.10.1-fr}
  \operatorname{Exalcotop}_{B/A}(C,L)\to
  \operatorname{Hom.\ cont}_C(\Upsilon_{C/B/A},L)
\end{equation}
compte tenu de \sref{0.18.5.3.1-fr} : pour tout idéal ouvert
$\mathfrak{M}$ de $A$, tout idéal ouvert $\mathfrak{N}$ de $B$ tel que
$\mathfrak{M}B\subset\mathfrak{N}$, tout idéal ouvert $\mathfrak{P}$ de $C$
tel que $\mathfrak{N}C\subset\mathfrak{P}$ et que $\mathfrak{P}.L=0$, toute
$(B/\mathfrak{N})$-extension $E$ de $C/\mathfrak{P}$ par $L$ qui est
$(A/\mathfrak{M})$-triviale provient de la donnée d'un $C$-homomorphisme
continu $f$ de $\Omega^1_{B/A}\otimes_B C$ dans $L$, et l'homomorphisme
\sref{0.20.7.10.1-fr} fait correspondre à l'image dans
$\operatorname{Exalcotop}_{B/A}(C,L)$ de la classe de $E$, la restriction de
$f$ à $\Upsilon_{C/B/A}$, dit \emph{homomorphisme caractéristique de $E$} et
noté encore $\chi_E$.
\end{env}

\begin{proposition}[20.7.11]
\label{0.20.7.11-fr}
Supposons que la topologie de $C$ soit telle que le carré d'un idéal ouvert
soit ouvert. Si $C$ est une $A$-algèbre topologique formellement lisse et si
$\Omega^1_{C/B}$ est un $C$-module formellement projectif, on a un
isomorphisme canonique
\begin{equation}
\tag{20.7.11.1}
\label{0.20.7.11.1-fr}
  \operatorname{Exalcotop}_B(C,L)\xrightarrow{\sim}
  \operatorname{Hom.\ cont}_C(\Upsilon_{C/B/A},L)
\end{equation}
pour tout $C$-module discret $L$ annulé par un idéal ouvert de $C$.
\end{proposition}

\begin{proof}
On a en effet, dans la suite exacte \sref{0.20.3.7.1-fr},
$\operatorname{Exalcotop}_A(C,L)=0$ \sref{0.19.4.4-fr}, donc
$\operatorname{Exalcotop}_B(C,L)=\operatorname{Exalcotop}_{B/A}(C,L)$ et
l'homomorphisme \sref{0.20.7.8.1-fr} n'est autre que
\sref{0.20.7.10.1-fr} ; le fait qu'il est bijectif se déduit de
\sref{0.20.6.13-fr} par passage à la limite inductive, compte tenu de ce que
la topologie de $\Omega^1_{C/B}$ est alors déduite de celle de $C$
\sref{0.20.4.5-fr} et de \sref{0.19.2.4-fr}.
\end{proof}

\begin{env}[20.7.12]
\label{0.20.7.12-fr}
Dans \sref{0.20.6.14-fr} on munit encore $\Upsilon^C_{B/A/\Lambda}$ de la
topologie induite par celle de $\Omega^1_{A/\Lambda}\otimes_A C$ et alors
l'homomorphisme \sref{0.20.6.14.6-fr} est continu, lorsque les anneaux
considérés sont topologiques et les homomorphismes d'anneaux continus.
\end{env}

\begin{env}[20.7.13]
\label{0.20.7.13-fr}
Si, dans \sref{0.20.6.23-fr}, on suppose que les anneaux $\Lambda$, $A$, $B$,
$C$ sont topologiques, les homomorphismes $s$, $u$, $v$ continus et $L$ un
$C$-module discret annulé par un idéal ouvert de $C$, alors on peut passer à
la limite inductive comme dans \sref{0.20.3.6-fr}, et on a une suite exacte
\begin{equation}
\tag{20.7.13.1}
\label{0.20.7.13.1-fr}
\begin{split}
  0\to\operatorname{D\acute{e}r.\ cont}_B(C,L)\to
  \operatorname{D\acute{e}r.\ cont}_A(C,L)\to
  \operatorname{D\acute{e}r.\ cont}_A(B,L)\to \\
  \to\operatorname{Exalcotop}_{B/\Lambda}(C,L)\to
  \operatorname{Exalcotop}_{A/\Lambda}(C,L)\to
  \operatorname{Exalcotop}_{A/\Lambda}(B,L)\to0.
\end{split}
\end{equation}
\end{env}

\begin{env}[20.7.14]
\label{0.20.7.14-fr}
Soient $A$ un anneau topologique, $B$ une $A$-algèbre topologique,
$\mathfrak{M}'$ un idéal ouvert de $A$, $\mathfrak{N}'$ un idéal ouvert de
$B$ tels que $\mathfrak{M}'B\subset\mathfrak{N}'$ ; posons
$A'=A/\mathfrak{M}'$,
\oldpage[0\textsubscript{IV-1}]{151}
$B'=B/\mathfrak{N}'$ ; le noyau de l'homomorphisme
$B\otimes_A B\to B'\otimes_{A'}B'$ est
$\mathfrak{U}'=\operatorname{Im}(\mathfrak{N}'\otimes B+B\otimes\mathfrak{N}')$,
d'où résulte aussitôt que le noyau de l'homomorphisme
\begin{equation}
\tag{20.7.14.1}
\label{0.20.7.14.1-fr}
  \varphi_{(\mathfrak{M}',\mathfrak{N}')}:
  \Omega^1_{B/A}\to\Omega^1_{B'/A'}
\end{equation}
est $((\mathfrak{J}\cap\mathfrak{U}')+\mathfrak{J}^2)/\mathfrak{J}^2$ ; par
ailleurs, l'homomorphisme \sref{0.20.7.14.1-fr} est surjectif, comme il
résulte de \sref{0.20.4.7-fr}. Si $\mathfrak{M}''$ (resp.\
$\mathfrak{N}''$) est un second idéal ouvert de $A$ (resp.\ $B$) tel que
$\mathfrak{M}''\subset\mathfrak{M}'$, $\mathfrak{N}''\subset\mathfrak{N}'$
et $\mathfrak{M}''B\subset\mathfrak{N}''$, et si l'on pose
$A''=A/\mathfrak{M}''$, $B''=B/\mathfrak{N}''$, on a de même un
homomorphisme surjectif
\[
  \varphi_{(\mathfrak{M}'',\mathfrak{N}''),(\mathfrak{M}',\mathfrak{N}')}:
  \Omega^1_{B''/A''}\to\Omega^1_{B'/A'}
\]
et ces homomorphismes forment évidemment un \emph{système projectif}. Si
l'on remarque que les
$((\mathfrak{J}\cap\mathfrak{U}')+\mathfrak{J}^2)/\mathfrak{J}^2$ forment
un système fondamental de voisinages de $0$ dans $\Omega^1_{B/A}$, on voit
donc que le \emph{séparé complété} $\widehat{\Omega}^1_{B/A}$ du $B$-module
topologique $\Omega^1_{B/A}$ est donné, à un isomorphisme canonique près, par
\begin{equation}
\tag{20.7.14.2}
\label{0.20.7.14.2-fr}
  \widehat{\Omega}^1_{B/A}=\varprojlim(\Omega^1_{B'/A'}).
\end{equation}

En outre, l'homomorphisme canonique
$j:\Omega^1_{B/A}\to\widehat{\Omega}^1_{B/A}$ est limite projective du
système projectif des $\varphi_{(\mathfrak{M}',\mathfrak{N}')}$, donc
$\varphi_{(\mathfrak{M}',\mathfrak{N}')}$ se factorise en
$\Omega^1_{B/A}\xrightarrow{j}\widehat{\Omega}^1_{B/A}\to
\Omega^1_{B'/A'}$, et comme il est surjectif on en conclut que
l'homomorphisme canonique
\begin{equation}
\tag{20.7.14.3}
\label{0.20.7.14.3-fr}
  \widehat{\Omega}^1_{B/A}\to\Omega^1_{B'/A'}
\end{equation}
est surjectif pour tout couple $(\mathfrak{M}',\mathfrak{N}')$.

On peut d'ailleurs dans ce qui précède remplacer partout $A'$ par $A$.

Enfin, si $L$ est un $B$-module topologique \emph{séparé et complet}, tout
$B$-homomorphisme continu de $\Omega^1_{B/A}$ dans $L$ se prolonge de façon
unique en un $\widehat B$-homomorphisme continu de
$\widehat{\Omega}^1_{B/A}$ dans $L$, et réciproquement un tel homomorphisme
redonne par composition avec $\Omega^1_{B/A}\to\widehat{\Omega}^1_{B/A}$ un
$B$-homomorphisme continu, si bien que l'on a un isomorphisme canonique
\[
  \operatorname{Hom.\ cont}_B(\Omega^1_{B/A},L)\xrightarrow{\sim}
  \operatorname{Hom.\ cont}_{\widehat B}(\widehat{\Omega}^1_{B/A},L).
\]
Plus particulièrement, si $L$ est un $B$-module discret annulé par un idéal
ouvert de $B$, on voit que l'isomorphisme canonique
\sref{0.20.4.8.2-fr} peut aussi s'écrire
\begin{equation}
\tag{20.7.14.4}
\label{0.20.7.14.4-fr}
  \operatorname{Hom.\ cont}_{\widehat B}(\widehat{\Omega}^1_{B/A},L)
  \xrightarrow{\sim}\operatorname{D\acute{e}r.\ cont}_A(B,L).
\end{equation}
\end{env}

\begin{proposition}[20.7.15]
\label{0.20.7.15-fr}
Soient $A$ un anneau discret, $B$ une $A$-algèbre topologique adique
($\mathbf{0}_{\mathrm I}$, 7.1.9), $\mathfrak{n}$ un idéal de définition de
$B$, $B_0=B/\mathfrak{n}$. On suppose que $\Omega^1_{B_0/A}$ et
$\mathfrak{n}/\mathfrak{n}^2$ sont des $B_0$-modules de type fini. Alors
$\widehat{\Omega}^1_{B/A}$ est un $\widehat B$-module de type fini.
\end{proposition}

\begin{proof}
Comme le carré de tout idéal ouvert de $B$ est ouvert, la topologie de
$\Omega^1_{B/A}$ est la topologie $\mathfrak{n}$-préadique
\sref{0.20.4.5-fr}. Compte tenu de l'hypothèse que $B$ est un anneau adique,
il suffit donc, en vertu de ($\mathbf{0}_{\mathrm I}$, 7.2.7 et 7.2.9), de
prouver
\oldpage[0\textsubscript{IV-1}]{152}
que $\Omega^1_{B/A}/\mathfrak{n}.\Omega^1_{B/A}=
\Omega^1_{B/A}\otimes_B B_0$ est un $B_0$-module de type fini. Mais cela
résulte de l'hypothèse et de la suite exacte \sref{0.20.5.12.1-fr}
\[
  \mathfrak{n}/\mathfrak{n}^2\to
  \Omega^1_{B/A}\otimes_B B_0\to\Omega^1_{B_0/A}\to0.
\]
\end{proof}

\begin{env}[20.7.16]
\label{0.20.7.16-fr}
La proposition \sref{0.20.7.15-fr} s'applique par exemple lorsque $A$ est un
corps $k$, $B=k'[[T_1,\ldots,T_n]]$ une algèbre de séries formelles munie de
sa topologie usuelle, $k'$ une extension \emph{finie} de $k$ (cf.\
\sref{0.21.9.2-fr}). On notera que le corps des fractions $K$ de $B$ a un
degré de transcendance infini sur $k$ ; lorsque $k$ est de caractéristique
$0$, on en déduit aussitôt (à l'aide de \sref{0.20.4.13-fr}, (i) et de
\sref{0.20.5.9-fr} notamment, en utilisant aussi le fait que toute dérivation
d'un corps de caractéristique $0$ se prolonge à toute extension) que
$\Omega^1_{B/k}$ \emph{n'est pas un $B$-module de type fini}.
\end{env}

\begin{env}[20.7.17]
\label{0.20.7.17-fr}
Soient $A$, $B$, $C$ trois anneaux topologiques, $u:A\to B$, $v:B\to C$ deux
homomorphismes continus ; remplaçant $A$, $B$, $C$ par des quotients par des
idéaux ouverts $A'=A/\mathfrak{M}'$, $B'=B/\mathfrak{N}'$,
$C'=C/\mathfrak{P}'$, avec $u(\mathfrak{M}')\subset\mathfrak{N}'$,
$v(\mathfrak{N}')\subset\mathfrak{P}'$, de sorte que l'on a des
homomorphismes $u':A'\to B'$, $v':B'\to C'$, on en déduit des homomorphismes
canoniques $u'_{C'/B'/A'}$, $v'_{C'/B'/A'}$ qui, en vertu de
\sref{0.20.5.4-fr}, forment des systèmes projectifs, et donnent par suite,
par passage à la limite, des homomorphismes canoniques, prolongements aux
séparés complétés des homomorphismes de la suite exacte
\sref{0.20.5.7.1-fr}
\[
  \Omega^1_{B/A}\otimes_B C\xrightarrow{v_{C/B/A}}\Omega^1_{C/A}
  \xrightarrow{u_{C/B/A}}\Omega^1_{C/B}\to0:
\]
\begin{align}
  \widehat v_{C/B/A}:\widehat{\Omega}^1_{B/A}
  \widehat\otimes_{\widehat B}\widehat C&\to\widehat{\Omega}^1_{C/A},
  \tag{20.7.17.1}\label{0.20.7.17.1-fr}\\
  \widehat u_{C/B/A}:\widehat{\Omega}^1_{C/A}&\to
  \widehat{\Omega}^1_{C/B}. \tag{20.7.17.2}\label{0.20.7.17.2-fr}
\end{align}
et dans la suite
\begin{equation}
\tag{20.7.17.3}
\label{0.20.7.17.3-fr}
  \widehat{\Omega}^1_{B/A}\widehat\otimes_{\widehat B}\widehat C
  \xrightarrow{\widehat v_{C/B/A}}\widehat{\Omega}^1_{C/A}
  \xrightarrow{\widehat u_{C/B/A}}\widehat{\Omega}^1_{C/B}\to0
\end{equation}
le composé de deux homomorphismes consécutifs est $0$, mais la suite
\emph{n'est pas nécessairement exacte}. Toutefois, si $B$ et $C$ sont
\emph{métrisables}, l'homomorphisme $\widehat u_{C/B/A}$ est \emph{surjectif},
et $\operatorname{Im}(\widehat v_{C/B/A})$ est \emph{dense} dans
$\operatorname{Ker}(\widehat u_{C/B/A})$ : cela résulte aussitôt
(cf.\ ($\mathbf{0}_{\mathrm I}$, 7.3.1)) de ce que, si
$(\mathfrak{M}_k)$ (resp.\ $(\mathfrak{N}_k)$) est une suite décroissante
d'idéaux de $B$ (resp.\ $C$), formant un système fondamental de voisinages de
$0$, et si l'on pose $B_k=B/\mathfrak{M}_k$, $C_k=C/\mathfrak{N}_k$, les
homomorphismes de transition
$\Omega^1_{C_{k+1}/B_{k+1}}\to\Omega^1_{C_k/B_k}$,
$\Omega^1_{C_{k+1}/A}\to\Omega^1_{C_k/A}$ et
$\Omega^1_{B_{k+1}/A}\to\Omega^1_{B_k/A}$ sont surjectifs
\sref{0.20.7.14-fr}.
\end{env}

\begin{proposition}[20.7.18]
\label{0.20.7.18-fr}
Soient $A$, $B$, $C$ trois anneaux topologiques, $u:A\to B$, $v:B\to C$
deux homomorphismes continus. On suppose $B$ et $C$ admissibles
($\mathbf{0}_{\mathrm I}$, 7.1.2) et métrisables. Pour que l'homomorphisme
canonique
\[
  \widehat v_{C/B/A}:\widehat{\Omega}^1_{B/A}
  \widehat\otimes_{\widehat B}\widehat C\to\widehat{\Omega}^1_{C/A}
\]
admette un inverse à gauche qui soit un $C$-homomorphisme continu, il faut et
il suffit que $C$ soit une $B$-algèbre formellement lisse relativement à $A$.
\end{proposition}

\begin{proof}
La condition est nécessaire en vertu de \sref{0.20.7.2-fr} et
\sref{0.19.1.6-fr}. Pour voir qu'elle est suffisante, notons que le $C$-module
topologique
$L=\widehat{\Omega}^1_{B/A}\widehat\otimes_{\widehat B}\widehat C$ est
métrisable et
\oldpage[0\textsubscript{IV-1}]{153}
complet en vertu de l'hypothèse, donc $E=D_C(L)$, muni de la topologie produit,
est métrisable et complet ; il est en outre admissible, car si $\mathfrak{R}$
est un idéal de définition dans $C$, la suite des
$(\mathfrak{R}\oplus L)^n=\mathfrak{R}^n\oplus\mathfrak{R}^{n-1}L$ tend vers
$0$. Comme l'application composée
\[
  D:B\xrightarrow{d_{B/A}}\Omega^1_{B/A}\to
  \widehat{\Omega}^1_{B/A}\widehat\otimes_{\widehat B}\widehat C=L
\]
est une $A$-dérivation continue de $B$ dans $L$, le $A$-homomorphisme continu
\[
  f:x\longmapsto(v(x),D(x))
\]
de $B$ dans $E$ définit sur $E$ une structure de $B$-extension. Comme $L$ est
un idéal fermé dans $E$, il résulte de \sref{0.19.9.5-fr} et de l'hypothèse
que l'application identique $C\to E/L$ (qui est un $B$-homomorphisme) se
factorise en $C\xrightarrow{g}E\to E/L$, où $g$ est un homomorphisme continu
tel que $g\circ v=f$ ; par suite \sref{0.20.1.3-fr}, $g$ est de la forme
$y\mapsto(y,D'(y))$, où $D'$ est une $B$-dérivation continue de $C$ dans $L$,
autrement dit $D'\circ v=D$. Compte tenu de \sref{0.20.7.14.4-fr}, on a
$D'=h\circ\widehat d_{C/A}$, où
$h:\widehat{\Omega}^1_{C/A}\to L$ est un $C$-homomorphisme continu ; mais on
a $\widehat d_{C/A}\circ v=\widehat v_{C/B/A}\circ D$ par définition, et
comme l'image de $B$ par $D$ engendre (topologiquement) le $C$-module $L$
\sref{0.20.4.7-fr}, la relation $D'\circ v=D$ donne bien
$h\circ\widehat v_{C/B/A}=1_L$.
\end{proof}

\begin{corollary}[20.7.19]
\label{0.20.7.19-fr}
Sous les hypothèses de \sref{0.20.7.18-fr}, si l'on suppose en outre que dans
$B$ et $C$ le carré d'un idéal ouvert soit ouvert, alors, pour que $C$ soit
une $B$-algèbre formellement lisse relativement à $A$, il faut et il suffit
que $\widehat v_{C/B/A}$ soit inversible à gauche.
\end{corollary}

\begin{proof}
En effet, les topologies de
$\widehat{\Omega}^1_{B/A}\widehat\otimes_{\widehat B}\widehat C$ et de
$\widehat{\Omega}^1_{C/A}$ sont alors déduites de celle de $C$
\sref{0.20.4.5-fr}, et tout $C$-homomorphisme de l'un dans l'autre est
nécessairement continu.
\end{proof}

\begin{env}[20.7.20]
\label{0.20.7.20-fr}
Soient $A$ un anneau topologique, $B$ une $A$-algèbre topologique
\emph{métrisable et complète}, $\mathfrak{R}$ un idéal \emph{fermé} de $B$,
$C=B/\mathfrak{R}$ l'anneau topologique quotient, qui est métrisable et
complet. Soit $(\mathfrak{M}_k)$ un système fondamental décroissant de
voisinages de $0$ dans $B$ formé d'idéaux, et posons
$B_k=B/\mathfrak{M}_k$,
$\mathfrak{R}_k=(\mathfrak{R}+\mathfrak{M}_k)/\mathfrak{M}_k$,
$C_k=B_k/\mathfrak{R}_k$. On a un système projectif d'homomorphismes
\[
  \delta_{C_k/B_k/A}:\mathfrak{R}_k/\mathfrak{R}_k^2\to
  \Omega^1_{B_k/A}\otimes_{B_k}C_k
\]
\sref{0.20.5.11.3-fr}, d'où l'on déduit en passant à la limite un
homomorphisme canonique
\[
  \widehat\delta_{C/B/A}:\mathfrak{R}/\mathfrak{R}^2\to
  \widehat{\Omega}^1_{B/A}\widehat\otimes_{\widehat B}\widehat C
\]
et en raisonnant comme dans \sref{0.20.7.17-fr}, on voit que l'homomorphisme
canonique
$\widehat u_{C/B/A}:\widehat{\Omega}^1_{B/A}
\widehat\otimes_{\widehat B}\widehat C\to\widehat{\Omega}^1_{C/A}$ est
\emph{surjectif} et que $\operatorname{Im}(\widehat\delta_{C/B/A})$ est
\emph{dense} dans $\operatorname{Ker}(\widehat u_{C/B/A})$.
\end{env}
